% !Mode:: "TeX:UTF-8"
% !TEX encoding = UTF-8 Unicode

%----------------------------------------------------------------------------------------
% Natural Language Processing: Neural Networks and Large Language Models
%
% Copyright 2021-2026
% Tong Xiao (xiaotong@mail.neu.edu.cn) and Jingbo Zhu (zhujingbo@mail.neu.edu.cn)
%----------------------------------------------------------------------------------------

%----------------------------------------------------------------------------------------
%    CONFIGURATIONS
%----------------------------------------------------------------------------------------

%\setcounter{chapter}{11}

%----------------------------------------------------------------------------------------
%	CHAPTER 12
%----------------------------------------------------------------------------------------

\chapter{推理能力（Reasoning）}

\thispagestyle{chapterurl} % add urls to the header

前面几章讨论了面向通用应用的大语言模型（large language model，LLM）：模型根据输入问题生成回答，能够自然地处理机器翻译、问答等多种任务。因此，许多自然语言处理（natural language processing，NLP）问题都可以表述为文本生成任务，使用大语言模型求解。在简单文本生成之外，一个自然的后续方向，是开发能够深入思考的模型，以解决更复杂的问题。尤其希望大语言模型能够生成中间思考步骤，从而表现出推理能力。

推理长期以来一直被视为人工智能（artificial intelligence，AI）最具挑战性的问题之一。NLP 社区曾尝试通过将句子分析为逻辑形式等方式处理这一问题，但这些工作主要依赖一系列流水线模块来构建推理模型，通常容易产生级联错误。相比之下，大语言模型为推理提供了更通用的框架：思维路径可以显式表示为文本，因此能够让模型在文本生成过程中进行推理。虽然推理能力可以视为大语言模型规模扩展所带来的涌现能力，但仍需要特定策略才能将其激活，包括思维链（chain of thought，CoT）提示、推理时扩展（inference-time scaling），以及训练大语言模型进行推理。

本章讨论提升大语言模型推理能力的方法。首先介绍与推理有关的基本概念，然后介绍开发推理模型的技术。由于前面章节已经详细讨论提示方法和思维链，本章重点关注两个研究方向：1）扩展推理时计算量；2）学习如何推理。最后，对全章作简要总结并讨论相关问题。

%----------------------------------------------------------------------------------------
% NEW SECTION
%----------------------------------------------------------------------------------------

\section{背景与概念}

这里定义推理问题，并讨论相关概念。

\subsection{什么是推理？}

人工智能研究历来将推理分为三种主要逻辑形式：\textbf{演绎推理}（deduction），即从已知前提推导具体结论，例如数学证明；\textbf{归纳推理}（induction），即从具体实例中归纳一般规则或模式；\textbf{溯因推理}（abduction），即为一组观察结果寻找最合理的解释~\cite{peirce1931:collected}。长期以来，主流机器学习和 NLP 研究主要关注归纳任务。例如，文本分类、命名实体识别等经典 NLP 应用都可以视为归纳任务：通过在大规模标注数据集上训练模型，使其学得能够泛化的映射函数。

然而，随着大语言模型发展，研究者正尝试解决复杂数学、智能体规划和科学发现等更困难的推理问题，这些问题往往需要严谨的演绎推理和溯因推理。现代大语言模型研究表明，标准预训练目标，即在海量文本语料上预测下一个词元，已经足够强大，能够帮助模型内化有关语言和世界的知识，其中也包括潜在的推理能力。不过，主要挑战在于，预训练阶段并未对推理作显式建模。因此，面对复杂问题时，大语言模型往往试图直接输出最终答案，而不是逐步展开逻辑推导。

例如，考虑一道简单的应用题：

\begin{quote}
\textit{如果 John 的年龄是 Mary 的两倍，而 Mary 再过 3 年就是 10 岁，那么 John 现在几岁？}
\end{quote}

基础大语言模型可能仅凭“两倍”和“10”这两个关键词的关联，给出一个错误猜测，例如 20。具有推理能力的模型则可能先显式生成逻辑链，再作答：

\begin{quote}
\textit{Mary 现在 $10 - 3 = 7$ 岁，因此 John 现在 $7 \times 2 = 14$ 岁。}
\end{quote}

要使用大语言模型解决复杂的演绎和溯因任务，就必须显式激活这些潜在推理能力，既可以通过思维链等提示技术，也可以通过强化学习等进阶后训练范式。

为了更好地理解推理问题，可以借助认知心理学的双过程理论来看待大语言模型的能力，即系统 1 思维和系统 2 思维~\cite{kahneman:2011thinking}。

\begin{itemize}
\item \vspace{0.5cm} \textbf{系统 1（直觉思维）：}在人类认知中，系统 1 自动且快速地运行，几乎不费力，甚至完全不费力。在大语言模型中，系统 1 思维对应模型一次运行就直接输出最终答案的范式。这种模式效率较高，通常足以完成简单问答等归纳任务。然而，面对复杂的多步问题时，系统 1 思维容易出现逻辑跳跃和幻觉。
\item \vspace{0.3cm} \textbf{系统 2（审慎思维）：}这种模式缓慢且有意识，需要逐步进行逻辑推导。对大语言模型而言，当模型被要求在推理阶段分配额外计算资源，显式生成中间推理步骤时，就启用了系统 2 思维。模型将思考过程表示为文本，分解复杂任务，验证中间结论，并在得出最终答案之前保持连贯的逻辑链。
\end{itemize}
\vspace{0.3cm}

大语言模型推理的核心目标是激活系统 2 能力。要理解这一点，必须超越传统符号逻辑的视角。在经典人工智能中，推理通常被定义为在结构化知识库上应用硬编码逻辑规则。但在大语言模型中，推理本身就是文本生成过程，可以定义为生成一系列连贯中间步骤、以得出最终答案的能力。在某种意义上，大语言模型的推理是一种为更难的问题动态分配更多推理计算量的机制。归根结底，大语言模型的推理能力，就是将潜在的系统 2 思维编码成一系列思考步骤，即推理步骤的能力。

\subsection{为什么推理能力很重要}

推理能力在 NLP 研究中很重要，因为它适用于以下几个富有挑战性的领域：

\begin{itemize}
\item \vspace{0.5cm} \textbf{数学：}求解数学问题需要严密逻辑。模型不能仅凭题目中的数字、关键词与相似示例答案之间的关联，就可靠地解决问题；它需要逐步应用数学规则，推导出正确结果。
\item \vspace{0.2cm} \textbf{编程：}编写软件远比简单预测下一个词更难。模型必须根据用户的具体要求和严格的代码逻辑生成代码。尤其在处理复杂问题时，将抽象要求转化为能够实现功能的代码，需要逐步推理。
\item \vspace{0.2cm} \textbf{智能体规划：}将大语言模型用作自主智能体时，它必须把庞大、复杂的目标分解成可执行的小步骤，还需要提前规划，并在出现问题时调整行动。
\item \vspace{0.2cm} \textbf{科学发现：}协助科学家需要考察数据、提出新假设，并认真检验。这高度依赖逐步的逻辑演绎和溯因推理。
\end{itemize}
\vspace{0.3cm}

除了求解更困难的问题，生成推理步骤也具有实际优势。最主要的是，它使模型具有可解释性。模型将思考过程写成文本，人们便能够阅读并检查其逻辑，甚至准确定位错误发生在哪一步。这使系统更容易验证，也更值得信任。此外，逐步推理有助于减少幻觉。面对难题时，模型不再随意猜测，而是以自身的逻辑步骤为最终答案提供依据，使输出更加可靠。

\subsection{从大语言模型到推理模型}

在大语言模型成为主流之前，人工智能中的推理主要依靠严格的规则和逻辑。研究者使用符号人工智能、语义分析等方法，将人类语言转换为所需形式，例如知识图谱或数据库查询。这些经典系统的逻辑虽然十分清楚、易于追踪，但系统本身很脆弱：一旦遇到不认识的词，或某条小规则失效，整个过程就会崩溃。

大语言模型出现后，推理问题被重新表述为文本生成任务，可以通过面向推理的技术更好地解决。大语言模型时代推理模型的发展，可以分为三个阶段：直接猜测答案、通过提示进行思考，以及最终学会如何思考。

早期语言模型主要以预测下一个词为训练目标。被问到问题时，它们试图立即输出答案。这种快速猜测适用于简单问答或写作任务，但面对复杂数学或逻辑问题时效果不佳。由于模型没有办法停下来，逐步梳理问题，因此经常犯错或产生幻觉。

第一次突破来自研究者对模型潜在推理能力的认识：大语言模型已经具备这些能力，只需要将其唤醒。只需使用思维链提示~\cite{wei-etal:2022chain}，要求模型“一步一步地思考”，就能迫使模型在给出答案之前先写出思考过程。不过，这种方法仍然高度依赖人工提示，并未真正教会模型逻辑背后的规则。

近期更具变革性的进展，是使用强化学习训练模型进行推理。强化学习不只是教模型生成文本，而是让模型通过试错来学习。模型找到正确的逻辑路径并得出正确答案时，就会获得奖励。随着训练推进，OpenAI 的 o1~\cite{openai:2024learning}、DeepSeek-R1~\cite{deepseek:2025r1} 等模型表现出了出色的推理能力，例如探索不同思路、反复检查自己的步骤，以及自行纠正错误。

这些自我纠错和探索能力催生了一个重要研究方向：推理时扩展 \cite{snell-etal:2024scaling}。这一思想很直接：在生成过程中分配更多计算资源，让模型更深入地思考，以解决更困难的问题。理想情况下，模型应根据问题的实际难度调整思考时间。不过，当前许多推理模型采用了更简单的方式——对所有问题都执行较长的思考过程。

从更广泛的角度看，增加推理阶段的计算量，普遍有助于提升推理能力。即使某个大语言模型未专门接受过推理训练，只需给予更多计算资源来“思考”，仍可增强其解决问题的能力。例如，可以让模型生成多条不同推理路径，通过投票选出最佳答案，也可以用搜索算法帮助模型找到正确的逻辑步骤。

推理模型也推动了近期基于大语言模型的智能体发展 \cite{yao-etal:2023react,wang-etal:2024survey}。AI 智能体若要自主完成任务，就需要借助推理，弄清如何一步一步地实现目标。现实任务很少能一帆风顺，因此智能体还必须意识到自己何时犯错，并知道如何纠正。这些都要求智能体具备良好的推理能力。

推理模型的快速发展，还带来了有趣但富有挑战性的研究方向 \cite{li-etal:2025system,chen-etal:2026towards}。一个方向关注效率。例如，如果模型学会动态调整思考时间，就不会在简单问题上浪费计算，而只会在困难问题上放慢速度。另一个方向是将推理从数学、编程等具有严格对错标准的领域，扩展到开放式现实情境。此外，使模型能够跨图像、音频、物理环境等不同类型的数据进行推理，也是构建更先进人工智能系统的关键一步。

%----------------------------------------------------------------------------------------
% NEW SECTION
%----------------------------------------------------------------------------------------

\section{通过提示激发推理}

提示是激发大语言模型推理能力的最简单方法。它不需要修改模型架构或参数，只需精心设计输入文本，就能促使模型输出推理过程。思维链提示是其中使用最广泛的方法之一，其核心思想是鼓励模型先生成中间逻辑推导，再给出最终答案。通常，可以在提示中提供几个逐步求解的示例，也可以直接在问题末尾添加“让我们一步一步地思考”这样的指令。下面给出一个思维链提示示例，其中突出显示了思维链指令。

\vspace{0.5em}
\begin{tcolorbox}[frame empty]

\begingroup
\setlength{\leftskip}{2em}
\setlength{\rightskip}{2em}

问：Tom 买了 3 双鞋，每双 $\$40$；又买了 2 件衬衫，每件 $\$20$。如果他在全部商品结账时使用一张 $10\%$ 的折扣券，总共需要支付多少钱？

\vspace{0.3em}

答：\ctext[RGB]{154, 245, 175}{让我们一步一步地思考。}

\vspace{0.3em}

\uline{首先，计算鞋子的总价：$3$ 双乘以 $\$40$，等于 $\$120$。接着，计算衬衫的总价：$2$ 件乘以 $\$20$，等于 $\$40$。折扣前总价为 $\$120 + \$40 = \$160$。折扣金额为 $10\%$ 乘以 $\$160$，即 $\$16$。最后，从总价中减去折扣金额：$\$160 - \$16 = \$144$。所以答案是 $\$144$。}

\endgroup

\end{tcolorbox}
\vspace{0.5em}

在这个例子中，大语言模型产生一条推理路径（以下划线标示），其中同时包含推理过程和最终答案。形式化地，给定思维链提示 $\mathbf{x}$\footnote{这里的思维链提示同时包括用户输入和思维链指令。}，模型通过最大化语言建模分数，生成最可能的输出，如下所示：
\begin{eqnarray}
\hat{\mathbf{y}} & = & \argmax_{\mathbf{y}} \Pr(\mathbf{y} | \mathbf{x}).
\end{eqnarray}

\noindent 这里，$\Pr(\mathbf{y} | \mathbf{x})$ 与标准语言模型中的文本生成概率相同。每个输出都是词元序列 $\mathbf{y}=y_0...y_n$，表示一条由若干推理步骤组成的推理路径：
\begin{eqnarray}
\mathbf{y} & = & \{\bar{\mathbf{y}}_1,...,\bar{\mathbf{y}}_{n_p}\}. \label{eq:reasoning-path-notation}
\end{eqnarray}

\noindent 其中，$n_p$ 为推理步骤数。每一步 $\bar{\mathbf{y}}_i$ 都是 $\mathbf{y}$ 的一个子序列，表示一次逻辑演绎或中间计算。需要注意，将推理路径划分为离散步骤有多种策略，例如按各个数学运算分段。为简单起见，在这个例子中将每个句子视为一个推理步骤，因此有

\[
\begin{array}{rl}
\bar{\mathbf{y}}_1: & \text{首先，计算鞋子的总价：}\nonumber \\
& \text{$3$ 双乘以 $\$40$，等于 $\$120$。} \nonumber \\
\bar{\mathbf{y}}_2: & \text{接着，计算衬衫的总价：} \nonumber \\
& \text{$2$ 件乘以 $\$20$，等于 $\$40$。} \nonumber \\
\bar{\mathbf{y}}_3: & \text{折扣前总价为 $\$120 + \$40 = \$160$。} \nonumber \\
\bar{\mathbf{y}}_4: & \text{折扣金额是 $10\%$ 乘以 $\$160$，即 $\$16$。} \nonumber \\
\bar{\mathbf{y}}_5: & \text{最后，从总价中减去折扣金额：$\$160 - \$16 = \$144$。} \nonumber \\
\bar{\mathbf{y}}_6: & \text{所以答案是 $\$144$。} \nonumber
\end{array}
\]

给定这样的推理路径，可以从最后一步提取答案。这种简单的思维链推理方法还有许多改进方式。例如，自一致性技术通过采样多条推理路径，并对最终答案进行多数投票，确定最可靠的答案 \cite{wang-etal:2023self}。其他策略包括：将复杂问题拆成易于处理的子任务的问题分解方法 \cite{zhou-etal:least}；提示模型评议并纠正自身逻辑错误的自我反思方法 \cite{shinn-etal:2023reflexion,madaan-etal:2024self}；以及高效探索多条分支推理路径的审慎推理框架 \cite{yao-etal:2024tree}。


前文已介绍多种基于思维链的提示技术，这里不再展开。读者可以回顾第 \ChapterLLMPrompting 章讨论的相关方法。


%----------------------------------------------------------------------------------------
% NEW SECTION
%----------------------------------------------------------------------------------------

\section{推理阶段的推理能力}
\label{sec:inference-time-reasoning}

推理阶段的推理能力与第 \ChapterLLMInference 章介绍的推理时扩展密切相关。本节在前面讨论的基础上，进一步关注如何提高大语言模型的推理表现。

\subsection{扩展推理时计算量}

大语言模型的成功，在很大程度上依赖计算资源的扩展。一种常见的改进方式是训练时扩展：只需使用更多训练数据、更大的模型规模和更多计算资源，研究者就构建出了十分强大的大语言模型 \cite{kaplan-etal:2020scaling,hoffmann-etal:2022training}。

然而，对于复杂推理任务，仅扩展训练过程还不够。推理需要逐步的逻辑推导和规划，与简单文本生成不同。训练时不可能向模型展示全部可能的推理路径。如果完全依赖训练来解决困难数学或逻辑问题，就需要极为庞大的数据量。理想情况下，希望模型能够利用简单逻辑步骤，即时构建自己的推理路径。因此，在模型推理阶段激活这种推理能力更有吸引力。

由此，一个自然的想法是在推理阶段增加计算量，让大语言模型在给出最终答案前，花更多时间探索不同推理步骤。这一方式的明显优势是，能够更灵活、更高效地使用计算资源。例如，对于复杂数学或逻辑问题，模型可以分配更多时间和算力，从而“思考”得更久。许多情况下，给予较小模型额外推理时间，甚至能使其超越被要求立即作答的大得多的模型。表 \ref{tab:scaling-comparison} 汇总比较了训练时扩展与推理时扩展。

\begin{table}[t]
\centering
\begingroup
\renewcommand{\arraystretch}{1.3}
\begin{tabular}{R{0.20\linewidth} | L{0.33\linewidth} | L{0.33\linewidth}}
特征 & 训练时扩展 & 推理时扩展 \\
\hline
工作方式 & 在初始学习阶段增大模型，并向其提供更多数据。 & 保持模型规模固定，但在作答前给予更多计算资源来“思考”，例如探索不同路径。 \\ \hline
训练成本 & 极高。成本昂贵，出问题后也难以修正。 & 低至中等。可能需要投入资源训练较小模型，用来检查答案。 \\ \hline
使用成本 & 可预测。生成一个词元的成本，主要取决于模型规模和上下文长度。 & 灵活，但可能很高。模型“思考”得越深，成本越高。\\ \hline
推理阶段硬件需求 & 需要较高的内存带宽，以快速读写数据。 & 同样需要较高带宽，还需要较大的内存容量，以探索多条长推理路径。 \\ \hline
数据需求 & 海量通用文本。 & 高质量、逐步求解的示例，例如数学证明，以及明确判断答案对错的评分方式。
\end{tabular}
\caption{训练时扩展与推理时扩展的比较}
\label{tab:scaling-comparison}
\endgroup
\end{table}

可以将推理时扩展视为一个\textbf{搜索扩展}（search scaling）问题。给予模型额外计算资源来“思考”，本质上就是允许模型探索更大的逻辑步骤空间，显著增加找到最优解的机会。这引出了三个关键问题：

\begin{itemize}
\item \vspace{0.5em} 如何生成不同的推理路径？
\item \vspace{0.3em} 如何在可能的步骤中搜索，找到最佳步骤？
\item \vspace{0.3em} 如何检查中间步骤和最终答案是否正确？
\end{itemize}
\vspace{0.5em}

下面几个小节围绕三个主题回答这些问题：采样多条推理路径（第~\ref{sec:sampling-reasoning-paths} 节）、推理的搜索策略（第~\ref{sec:search-strategies-for-reasoning} 节），以及验证（第~\ref{sec:reasoning-verification} 节）。

\subsection{采样多条推理路径}
\label{sec:sampling-reasoning-paths}

扩展推理时计算量最直接的方式，是针对单个输入生成多条推理路径，而不只生成一条。通过生成多条不同路径，大语言模型有更多机会找到正确逻辑及相应答案。

为此，通常在文本生成过程中引入随机性。一种典型方法是使用温度采样（temperature sampling），并结合 top-$k$ 采样或核采样，即 top-$p$ 采样。在标准自回归生成中，对于每个位置 $i$，大语言模型输出概率分布 $\Pr(\cdot | \mathbf{x}, \mathbf{y}_{< i})$；该分布定义在词表 $V$ 上，用于预测下一个词元 $y_i$，依据对应的未归一化分数 $u_{y_i}$（logit）计算。温度采样在应用 Softmax 函数之前，用温度参数 $T > 0$ 缩放这些分数，修改该分布：
\begin{eqnarray}
\Pr(y_i | \mathbf{x}, \mathbf{y}_{< i}) & = & \frac{\exp(u_{y_i} / T)}{\sum_{y'_i \in V} \exp(u_{y'_i} / T)}.
\end{eqnarray}

\noindent 较高温度会使分布更平坦，鼓励模型探索更多样的推理步骤。然而，不加约束的随机性可能偶尔从低概率词元的“长尾”中采样，生成退化或不合逻辑的文本。为保持生成路径的质量，温度采样通常与词表截断方法结合：

\begin{itemize}
\item \vspace{0.5em} \textbf{Top-$k$ 采样}：将采样池限制为概率最高的 $k$ 个后续词元，并在它们之间重新分配概率质量。
\item \vspace{0.3em} \textbf{核采样（top-$p$ 采样）}：采用动态截断策略，选择排名靠前、累积概率超过阈值 $p$ 的最小词元集合。
\end{itemize}
\vspace{0.5em}

采用这些随机解码策略后，给定输入 $\mathbf{x}$，大语言模型可以采样得到一组多样化输出序列，即候选推理路径：
\begin{eqnarray}
\mathcal{Y}^{\mathrm{top}} & = & \{\hat{\mathbf{y}}^1, \dots, \hat{\mathbf{y}}^N\}.
\end{eqnarray}

\noindent 其中，$N$ 表示推理时计算预算。$N$ 越大，探索不同推理路径所用的计算量就越多。得到这一组多样化路径后，接下来的问题是如何选出最佳路径。为此，常用两种方法：$N$ 次采样择优和自一致性。

\subsubsection{$N$ 次采样择优（BoN）}
第 \ChapterLLMTuning\ 章已介绍 $N$ 次采样择优（Best-of-N，BoN），将其作为推理时对齐的一种方法。在那里，使用奖励模型 $r(\mathbf{x}, \hat{\mathbf{y}})$，选出最符合人类偏好的输出。

对推理任务而言，BoN 采样的概念相同，但目标有所变化。“评判者”通常是验证器（verifier），而非偏好奖励模型，可以将其记为通用评分函数 $v$，用于评估推理路径在逻辑上是否合理。验证器为每条采样路径赋予分数 $v(\mathbf{x}, \hat{\mathbf{y}})$，再选择使分数最大的最佳推理路径 $\hat{\mathbf{y}}_{\mathrm{best}}$：
\begin{eqnarray}
\hat{\mathbf{y}}_{\mathrm{best}} & = & \argmax_{\hat{\mathbf{y}} \in \{\hat{\mathbf{y}}^1, \dots, \hat{\mathbf{y}}^N\}} v(\mathbf{x}, \hat{\mathbf{y}}).
\end{eqnarray}

\noindent 然后，从得分最高的路径中提取最终答案。可以用一个直观类比理解 BoN 过程：设想教室里有 $N$ 个学生正在解一道数学题。老师（验证器）给所有人的草稿打分，再选出得分最高的学生。

BoN 采样虽然有效，但效果依赖验证器的质量。如果验证器能力不足，就很容易被有缺陷、却碰巧得出正确最终答案的逻辑欺骗，这就是奖励投机（reward hacking）。验证器的设计将在第~\ref{sec:reasoning-verification} 节详细讨论。

\subsubsection{自一致性}
除了用学习得到的验证器选择最佳路径，还可以采用自一致性。它所依据的直觉是：在复杂推理任务中，往往存在多条不同逻辑路径通向同一个正确答案，而错误推理路径可能各不相同，产生随机错误。自一致性不为推理过程评分，而是直接通过对最终答案进行多数投票，将推理路径边缘化。形式化地，令 $A$ 为所有可能最终答案的集合。对于每条采样路径 $\hat{\mathbf{y}}^k$，提取其最终预测答案 $\hat{a}^k = \mathrm{Extract}(\hat{\mathbf{y}}^k)$。自一致性策略选择答案 $\hat{a}_{\mathrm{sc}}$，使其在 $N$ 条采样路径中出现最频繁：
\begin{eqnarray}
\hat{a}_{\mathrm{sc}} & = & \argmax_{a \in A} \sum_{k=1}^N \mathbbm{1}(\mathrm{Extract}(\hat{\mathbf{y}}^k) = a),
\end{eqnarray}

\noindent 其中，$\mathbbm{1}(\cdot)$ 为指示函数。

\subsubsection{平衡采样的多样性与质量}

在通过采样多条推理路径来扩展推理时计算量时，需要在多样性与质量之间权衡。推理任务通常希望采样路径足够多样，以增加其中包含正确逻辑的机会。一个常用启发式方法是提高采样温度。分布变得更平坦后，采样就不会只集中在少数高概率路径上。因此，低概率路径更容易被选中，模型能够在 $N$ 个样本中探索更多不同推理路径。不过，随着温度升高，样本的多样性虽然增加，单条路径的平均逻辑质量却可能下降。

为平衡多样性与质量，一种直接的方法是在推理阶段引入显式惩罚项或奖励项，修改搜索目标。在这种设定下，可以将多样性建模为一个函数，与原生成概率相结合。形式化地，修改后的搜索目标为
\begin{eqnarray}
\mathrm{score}(\mathbf{x},\mathbf{y}) & = & \log \Pr(\mathbf{y}|\mathbf{x}) + \lambda D(\mathbf{y},\mathcal{Y}^{\mathrm{top}}). \label{diversity-quality-inference}
\end{eqnarray}

\noindent 其中，$\log \Pr(\mathbf{y} | \mathbf{x})$ 倾向于概率更高的输出，而 $D(\mathbf{y},\mathcal{Y}^{\mathrm{top}})$ 是惩罚项或奖励项，用来鼓励相对于已有采样路径 $\mathcal{Y}^{\mathrm{top}}$ 的多样性。超参数 $\lambda$ 控制这一权衡的强度。
$D(\cdot)$ 的设计可以十分灵活。例如，一种简单惩罚可以降低词元或 $n$ 元语法的分数，前提是它们已在 $\mathcal{Y}^{\mathrm{top}}$ 中频繁出现。反过来，也可以设计奖励函数，奖励显著的语义差异。这样，大语言模型就能够产生新的推理结构，而不只是对同一思路作细微改写。

从更广泛的角度看，这种采样困境也与强化学习中的探索—利用权衡有关 \cite{Sutton-and-Barto:2018RL}。在基于强化学习的调优中，大语言模型可以视为策略 $\pi_\theta$，生成推理路径 $\mathbf{y}$，用于回答问题 $\mathbf{x}$。一个通用的正则化目标可以写为
\begin{equation}
\max_{\pi_\theta} \mathbb{E}_{\mathbf{y}\sim\pi_\theta} \left[R(\mathbf{x},\mathbf{y})\right] - \beta D_{\mathrm{KL}} \left( \pi_\theta(\cdot|\mathbf{x}) \,\|\,  \pi_{\mathrm{ref}}(\cdot|\mathbf{x}) \right) + \alpha \mathrm{Entropy} \left( \pi_\theta(\cdot|\mathbf{x}) \right),
\label{diversity-quality-rl}
\end{equation}

\noindent 其中，$R(\mathbf{x},\mathbf{y})$ 衡量生成推理路径的质量。KL 散度项约束学习到的策略，使其保持接近参考策略，主要目的是防止策略过度偏移，从而减少奖励投机或语言质量下降等不良行为。熵项则显式鼓励策略保持不那么集中的分布，促进训练中的探索。

强化学习目标与推理阶段的多样性—质量权衡之间存在概念联系，尽管两者作用于不同空间。推理时，搜索过程控制在固定模型生成的推理路径中探索的广度；在基于强化学习的调优中，KL 正则化控制学习到的策略可以偏离参考策略多远。较大的 KL 系数使策略保持接近参考模型，较小系数则允许优化过程探索更广泛的策略空间，寻找奖励更高的行为。这类策略变化可能使模型生成在参考策略下概率很低的推理路径，从而产生更丰富的输出。

推理时搜索和采用 KL 正则化的强化学习，都涉及保守行为与更广泛探索之间的平衡。区别在于，推理时方法在固定模型下探索不同输出，而强化学习通过更新模型参数搜索策略空间。两者中，更广泛的探索都可能发现更好、更多样的推理行为，但过度探索可能降低质量或稳定性。

\subsection{推理的搜索策略}
\label{sec:search-strategies-for-reasoning}

在实践中，寻找最优推理路径可以直接表述为经典搜索问题：从初始状态出发，采取动作生成一个推理步骤，将当前状态扩展为一组新状态，反复进行，直到满足终止条件。第 \ChapterLLMInference 章已经介绍了树搜索、束搜索等多种经典算法，这里不再赘述算法细节，而是重点讨论如何将推理路径搜索映射到这些经典搜索框架中。

要将搜索算法用于大语言模型推理，需要定义搜索问题的标准组成部分：

\begin{itemize}
\item \vspace{0.5em} \textbf{搜索状态}：在大语言模型推理中，状态就是模型截至当前构建的上下文。具体而言，第 $t$ 步的状态 $s_t$ 由原始问题 $\mathbf{x}$ 和截至当前生成的部分推理路径共同构成：$s_t = (\mathbf{x}, \bar{\mathbf{y}}_1, \dots, \bar{\mathbf{y}}_t)$。
\item \vspace{0.3em} \textbf{动作}：生成下一个推理步骤，即文本片段 $\bar{\mathbf{y}}_{t+1}$。
\item \vspace{0.3em} \textbf{转移模型}：执行动作只是将新文本片段追加到已有状态，产生下一状态：$s_{t+1} = (s_t, \bar{\mathbf{y}}_{t+1}) = (\mathbf{x}, \bar{\mathbf{y}}_1, \dots, \bar{\mathbf{y}}_t, \bar{\mathbf{y}}_{t+1})$。
\item \vspace{0.3em} \textbf{状态评估函数}：对当前状态通向最终答案的逻辑正确性进行评分的函数。例如，可以使用硬编码的启发式规则或训练得到的神经验证器来评估搜索状态。
\item \vspace{0.3em} \textbf{终止条件}：停止搜索的准则。当模型显式生成最终答案序列，或搜索达到预定义计算预算，例如已探索节点的最大深度时，搜索停止。
\end{itemize}
\vspace{0.5em}

将推理任务表述为搜索问题后，就可以使用各种搜索算法，高效探索推理空间。例如，可以用广度优先搜索、深度优先搜索等简单树搜索方法对大语言模型的推理过程建模。模型在每一步探索多个候选思路并进行评估，决定是否沿某条路径继续前进，或在遇到死路时回溯。对于不确定性更高或奖励延迟的任务，常使用蒙特卡洛树搜索。这种方法帮助模型平衡对新推理路径的探索与对已知步骤的利用，有时还会结合外部环境反馈和自我反思来改进搜索 \cite{yao-etal:2024tree,hao-etal:2023reasoning,zhou-etal:2024language}。

在树结构中，多条不同推理路径可能汇聚到同一个中间结论，从而产生冗余。为解决这一问题，一些方法将推理空间建模为图，例如有向无环图（directed acyclic graph，DAG）\cite{zhang-etal:2025cumulative}。这样，模型就能够构建更高效的搜索空间结构，捕捉逻辑演绎中非线性的相互依赖。除了使用外部搜索算法，还可以将搜索动态直接内化到大语言模型参数中。例如，通过在 A* 等启发式搜索算法的执行轨迹上训练模型，可以使其学会隐式生成搜索过程和最优规划，而不依赖外部控制脚本 \cite{lehnert-etal:2025beyond}。

通过搜索扩展推理时计算量，是一种强大的策略，但搜索投入与推理表现之间并非简单的线性关系。直觉上，也许会认为，给予模型更多探索时间，就能持续得到更好的答案；但实际上，性能收益很快就会进入平台期。一个主要原因是模型固有的能力上限：如果它缺乏解决问题所需的知识，再多搜索也无法凭空产生正确逻辑。此外，随着搜索投入增加，生成路径的多样性会下降，模型因而在缺乏新价值的冗余思路上浪费计算。

\subsection{验证}
\label{sec:reasoning-verification}

接下来讨论推理中的验证问题。虽然本节重点是推理阶段的方法，但推理模型所使用的许多验证器仍需训练，因此下面也会介绍训练这些模型的一些细节。

\subsubsection{过程奖励模型}

验证推理结果主要有两种范式：\mindex{结果奖励模型}（outcome reward models，ORM）和\mindex{过程奖励模型}（process reward models，PRM）\cite{uesato-etal:2022solving,lightman-etal:2024lets}。结果奖励模型仅评估推理路径的最终答案。虽然训练起来更容易，但它容易出现“假阳性”，因此往往不足以处理推理任务。考虑下面这道简单数学题：

\vspace{0.5em}
\begin{quote}
\textit{求解 $x$：$2x + 4 = 10$。}
\end{quote}

\vspace{0.5em}
模型可能生成如下推理路径：
\vspace{0.5em}

\begin{quote}
\textit{$2x + 4 = 10$ 意味着 $x$ 必须是 $10$ 的一半，即 $5$。}

\textit{但 $5 - 2 = 3$。}

\textit{因此，$x = 3$。}
\end{quote}
\vspace{0.5em}

虽然最终答案 $3$ 是正确的，但逻辑步骤完全错误。如果只使用结果奖励模型，这条路径可能得到高分，无意间强化错误逻辑。为此，可以使用过程奖励模型评估推理路径中每一步的正确性。给定推理路径 $\mathbf{y} = \{\bar{\mathbf{y}}_1, \dots, \bar{\mathbf{y}}_{n_p}\}$，过程奖励模型为当前步骤赋予正确性分数 $r_t$，评分对象为 $\bar{\mathbf{y}}_t$，条件是输入 $\mathbf{x}$ 和此前所有步骤 $\bar{\mathbf{y}}_{<t}$。可以将步骤级验证函数定义为
\begin{eqnarray}
r_t & = & v_{\mathrm{prm}}(\bar{\mathbf{y}}_t | \mathbf{x}, \bar{\mathbf{y}}_{<t}).
\end{eqnarray}

\noindent 其中，$v_{\mathrm{prm}}(\cdot)$ 是一个经过训练、用于预测正确性分数的神经网络。例如，可以在现有大语言模型的输出表示上附加一个简单分类头，如线性层后接 sigmoid 激活函数。

为确定整条推理路径的总体质量，可以汇总这些步骤级分数。一种简单做法是取所有步骤分数的最小值。换言之，只有当推理的“最薄弱环节”在逻辑上也合理时，才认为整条推理正确。验证分数为
\begin{eqnarray}
v(\mathbf{x}, \mathbf{y}) & = & \min_{1 \le t \le n_p} v_{\mathrm{prm}}(\bar{\mathbf{y}}_t | \mathbf{x}, \bar{\mathbf{y}}_{<t}).
\end{eqnarray}

\noindent 也可以将验证分数定义为各步骤分数的平均值：
\begin{eqnarray}
v(\mathbf{x}, \mathbf{y}) & = & \frac{1}{n_p}\sum_{1 \le t \le n_p} v_{\mathrm{prm}}(\bar{\mathbf{y}}_t | \mathbf{x}, \bar{\mathbf{y}}_{<t}).
\end{eqnarray}

训练验证模型需要标注数据集。结果奖励模型的训练只需知道路径末尾的正确答案，而过程奖励模型需要步骤级监督。构建过程奖励模型训练数据集，通常有两种方式：

\begin{itemize}
\item \vspace{0.5em} \textbf{人工标注：}人类专家手动求解问题，将大语言模型生成的每个中间步骤标为“正确”或“错误”。这种方法能够训练出准确的验证器，但成本高，难以扩展；对于需要领域专家的高等数学或编程任务，尤其如此。
\item \vspace{0.3em} \textbf{通过轨迹采样自动验证：}还可以用自动化方法估计一个步骤的价值 \cite{wang-etal:2024math}。例如，如果拥有已知最终答案的数据集，如数学竞赛数据，就可以采用\mindex{蒙特卡洛轨迹采样}（Monte Carlo rollouts）。从中间步骤 $\bar{\mathbf{y}}_t$ 出发，提示大语言模型继续生成多条完整路径。如果这些路径中有较高比例最终得到正确标准答案 $a^*$，就可以认为中间步骤 $\bar{\mathbf{y}}_t$ 处在正确方向上。形式化地，一个步骤的估计价值 $\hat{r}_t$ 可以定义为到达正确答案的期望概率：
    \begin{eqnarray}
    \hat{r}_t & = & \mathbb{E}_{\mathbf{y}_{>t} \sim \pi_\theta} [ \mathbbm{1}(\mathrm{Extract}(\mathbf{y}_{>t}) = a^*) ].
    \end{eqnarray}
    \noindent 其中，$\mathbf{y}_{>t}$ 是从 $\bar{\mathbf{y}}_t$ 开始的剩余路径，$\mathrm{Extract}(\mathbf{y}_{>t})$ 表示从路径 $\mathbf{y}_{>t}$ 中提取的答案。通过轨迹采样估计这些值后，就可以训练过程奖励模型预测 $\hat{r}_t$。
\end{itemize}\vspace{0.5em}


\subsubsection{生成式验证器}

上面所有验证器都需要额外的架构修改和训练工作，例如训练分类头，输出数值分数。鉴于大语言模型具有很强的文本生成能力，另一种自然方式是直接提示模型评估推理步骤，并以自然语言生成验证反馈。这称为\textbf{生成式验证}（generative verification）\cite{zheng-etal:2023judging}。

这种方式将验证任务表述为标准文本生成问题，而不是修改模型架构以输出单个标量。将问题和中间步骤重新输入大语言模型，就可以让它充当评判者。此时，模型就像一位在试卷页边写评语的老师，而不只是给出最终成绩。

下面通过一个简单例子说明其工作方式。假设模型正在求解一个方程组：
\vspace{0.5em}
\begin{quote}
\textit{题目：求解 $x$ 和 $y$，已知 $2x + y = 10$ 和 $x - y = 2$。}
\end{quote}
\vspace{0.5em}

假设模型已经生成两个正确步骤：
\vspace{0.5em}
\begin{quote}
\textit{由第二个方程，可以将 $x$ 表示为 $y$ 的形式：$x = y + 2$。}

\textit{将 $x$ 代入第一个方程：$2(y + 2) + y = 10$。}
\end{quote}
\vspace{0.5em}

现在，模型试图生成第三步，但其中包含计算错误：

\vspace{0.5em}
\begin{quote}
\textit{展开方程得到 $2y + 2 + y = 10$，化简为 $3y + 2 = 10$。}
\end{quote}
\vspace{0.5em}

可以将原问题、步骤 1 和步骤 2、当前步骤 3，与验证提示组合起来，例如“\textit{请根据前面的步骤，检查当前步骤的逻辑是否正确，并解释你的结论。}”生成式验证器随后可能输出如下文本评议：
\vspace{0.5em}
\begin{quote}
\textit{评议：Incorrect（错误）。在当前步骤中，将步骤 2 中的 $2$ 分配到 $(y + 2)$ 时，必须将括号内两项都乘以 $2$。结果应为 $2y + 4$，而非 $2y + 2$。正确展开式为 $2y + 4 + y = 10$，化简为 $3y + 4 = 10$。}
\end{quote}
\vspace{0.5em}

为形式化描述这一生成式验证过程，假设原始输入 $\mathbf{x}$ 由问题 $\mathbf{z}$ 和指令 $\mathbf{c}$ 构成，指令要求模型根据上下文验证当前步骤。可以将问题 $\mathbf{z}$、推理历史 $\bar{\mathbf{y}}_{\le t}$ 和指令 $\mathbf{c}$ 拼接成一段长文本，再将其输入大语言模型，得到自然语言评议 $\mathbf{p}_t$：
\begin{eqnarray}
\mathbf{p}_t & \sim & \Pr(\cdot | \mathbf{z}, \bar{\mathbf{y}}_{\le t}, \mathbf{c}).
\end{eqnarray}

\noindent 其中，$\Pr(\cdot)$ 表示现有大语言模型所建模的概率分布。因此，无需训练，就可以直接使用强大的大语言模型进行准确验证。图 \ref{fig:pre-generative-verfication} 比较了这一生成式方法与传统判别式验证器。

\begin{figure}[!t]
\centering
% !Mode:: "TeX:UTF-8"
% !TEX encoding = UTF-8 Unicode

\begin{tikzpicture}
	
\def\ssep{0.3cm}

\tikzstyle{lnode} = [minimum width=5.0cm,minimum height=1.7cm,draw,thick,inner sep=2pt]

%%% scalar verifier
\begin{scope}

\node [lnode,anchor=west] (llm) at (0,0) {};
\node [anchor=south] (llmlabel) at ([yshift=0.4cm]llm.south) {\large{LLM}};

\node [anchor=center] (z) at ([xshift=0.3cm,yshift=-0.9cm]llm.south west) {$\mathbf{z}$};
\node [anchor=center] (steps) at ([xshift=0.9cm]z.center) {$\bar{\mathbf{y}}_{\le t}$};
%\node [anchor=center] (c) at ([xshift=0.9cm]steps.center) {$\mathbf{c}$};

\draw [->,thick] ([yshift=0.3cm]z.center) -- ([yshift=0.8cm]z.center);
\draw [->,thick] ([yshift=0.3cm]steps.center) -- ([yshift=0.8cm]steps.center);
%\draw [->,thick] ([yshift=0.3cm]c.center) -- ([yshift=0.8cm]c.center);

\node [anchor=north east,draw,dashed,minimum width=2.5cm,inner sep=2pt] (hidden) at ([xshift=-0.2cm,yshift=-0.2cm]llm.north east) {\scriptsize{最后隐层}};

\node [anchor=south,draw,thick,minimum width=2.5cm] (classification) at ([yshift=0.9cm]hidden.north) {\scriptsize{分类头}};
\draw [->,thick] ([yshift=0.05cm]hidden.north) -- ([yshift=-0.1cm]classification.south);

\node [anchor=south,align=left] (prediction) at ([yshift=0.7cm]classification.north) {\footnotesize{以下结果的概率}\\[-0.1cm] \scriptsize{``Correct''}};
\draw [->,thick] ([yshift=0.1cm]classification.north) -- (prediction.south);

\node [anchor=north] (caption) at ([yshift=-1.4cm]llm.south) {(a) 判别式验证器};

\end{scope}

%%% generative verifier
\begin{scope}[xshift=7.5cm]

\node [lnode,anchor=west] (llm) at (0,0) {};
\node [anchor=south] (llmlabel) at ([yshift=0.4cm]llm.south) {\large{LLM}};

\node [anchor=center] (z) at ([xshift=0.3cm,yshift=-0.9cm]llm.south west) {$\mathbf{z}$};
\node [anchor=center] (steps) at ([xshift=0.9cm]z.center) {$\bar{\mathbf{y}}_{\le t}$};
\node [anchor=center] (c) at ([xshift=0.9cm]steps.center) {$\mathbf{c}$};

\draw [->,thick] ([yshift=0.3cm]z.center) -- ([yshift=0.8cm]z.center);
\draw [->,thick] ([yshift=0.3cm]steps.center) -- ([yshift=0.8cm]steps.center);
\draw [->,thick] ([yshift=0.3cm]c.center) -- ([yshift=0.8cm]c.center);

\node [anchor=south] (critique) at ([xshift=-1.5cm,yshift=0.7cm]llm.north east) {\footnotesize{文本评议}};

\draw [->,thick] ([yshift=-0.6cm]critique.south) -- ([yshift=-0.0cm]critique.south);


\node [anchor=north] (caption) at ([yshift=-1.4cm]llm.south) {(b) 生成式验证器};


\end{scope}
	
\end{tikzpicture}

\caption{判别式与生成式验证方法的比较。子图 (a) 展示判别式验证器：在大语言模型最后一个隐层上附加分类头，输出表示正确性的数值概率。子图 (b) 展示生成式验证器：将问题 $\mathbf{z}$、推理历史 $\bar{\mathbf{y}}_{\le t}$ 和指令 $\mathbf{c}$ 作为提示输入模型，生成自然语言评议，而非标量值。}
\label{fig:pre-generative-verfication}
\end{figure}

自然语言评议虽然有助于人类理解，但树搜索等部分搜索算法可能需要数值分数，才能对不同推理路径排序。从生成式验证器提取连续分数的一种常用方法，是衡量模型以指示词元（如“Correct”（正确）或“Incorrect”（错误））作为评议第一个词的概率。因此，数值验证分数可定义为：
\begin{eqnarray}
v(\mathbf{z}, \bar{\mathbf{y}}_{\le t}, \mathbf{c}) & = & \Pr(y_{\mathrm{anchor}} = \text{``Correct''} | \mathbf{z}, \bar{\mathbf{y}}_{\le t}, \mathbf{c}).
\end{eqnarray}

\noindent 其中，$y_{\mathrm{anchor}}$ 是指示词元，通常通过启发式规则从输出文本中提取。

相对于传统标量奖励模型，生成式验证器有几个优势。第一，它为问题求解提供可解释性。验证器用文本解释推理，开发者便能诊断一条推理路径为何被接受或拒绝。第二，它支持自我纠错。文本评议往往明确指出如何修正错误，例如上例中给出了正确展开式 $2y + 4$，因此推理系统可以直接吸收反馈，纠正当前状态并继续搜索，而不必简单丢弃整条路径。第三，它对架构的修改要求很小。只要现有大语言模型具有较强推理能力，同一个模型就可以同时用作问题求解器和验证器。

\subsubsection{用作验证器的形式符号系统}

神经验证器虽然应用广泛，但具有一个共同弱点：它们仍然是神经网络，因此也可能产生幻觉或犯逻辑错误。对于复杂数学、软件工程等要求严格精确性的任务，还可以使用形式符号系统（formal symbolic system）作为验证器。例如，在代码优化任务中，当模型尝试优化 LLVM-IR 等中间表示代码时，Alive2 等形式化等价性检查器可以自动验证，新生成的代码是否与原代码保持完全相同的语义。

形式系统严格遵循预定义的数学和逻辑规则，因此其判断是 100\% 可靠的。要将形式系统用作验证器，大语言模型必须将推理步骤生成为可执行代码或形式化数学语句，而不只是普通自然语言。例如，求解数学问题时，模型可以编写一小段 Python 脚本，执行某个中间计算，再由系统通过外部工具运行代码。验证过程十分直接：

\begin{itemize}
\item \vspace{0.5em} 如果代码运行无错误，并产生有效输出，就将这一推理步骤验证为正确。
\item \vspace{0.3em} 如果代码抛出错误，例如语法错误或违反数学规则，就将这一推理步骤验证为错误。
\end{itemize}
\vspace{0.5em}

形式化地，令中间步骤 $\bar{\mathbf{y}}_t$ 为生成的一段代码或形式逻辑，令 $E$ 表示形式执行环境，则验证分数可定义为：
\begin{eqnarray}
v(\mathbf{x}, \bar{\mathbf{y}}_{<t}, \bar{\mathbf{y}}_t) & = & \begin{cases}
1, & \text{若 } E(\bar{\mathbf{y}}_t) \text{ 返回正确结果，} \\
0, & \text{否则。}
\end{cases}
\end{eqnarray}


这一方法有一个明显优势：消除了验证器幻觉导致的“假阳性”。此外，执行失败时，形式系统通常会生成错误信息。将这些信息反馈给大语言模型，就能帮助它理解问题所在，在下一次尝试中纠正逻辑。

不过，使用形式符号系统也存在局限。第一，它要求大语言模型具备较强的编程或形式化能力，因为将抽象推理转化为严格代码并不容易。第二，形式系统虽然非常适合数学和编程，但不适合写文章、规划旅行等开放式现实任务，因为这类任务的推理难以映射为严格的逻辑规则。

%----------------------------------------------------------------------------------------
% NEW SECTION
%----------------------------------------------------------------------------------------

\section{学习如何推理}

提升大语言模型的推理能力很有挑战，通常需要结合多种方法。如前一节所述，推理时扩展旨在生成过程中激活现有模型的“思考”过程。另一种自然思路，是在训练时将推理能力直接内化到模型参数中。本节重点讨论这种“学习如何推理”的范式。

推理数据可以加入任何训练阶段。例如，预训练时可以直接将推理路径当作普通文本，进行标准的下一词元预测；后训练时也可以开展基于推理的训练。例如，在监督微调阶段，可以使用高质量的逐步解题样例微调模型，通常通过包含思维链轨迹的数据集来训练。此外，还可以利用强化学习进一步优化这些推理过程。这里重点讨论如何在大语言模型对齐中加入基于推理的监督。下面几个小节将讨论监督推理数据、过程监督、强化学习技术，以及扩展推理能力的策略。

\subsection{监督推理数据}

形式化地，在监督微调中，使用输入输出对数据集 $\mathcal{D} = \{(\mathbf{x}, \mathbf{y})\}$ 训练大语言模型，标准训练目标为
\begin{eqnarray}
\hat{\theta} & = & \argmax_{\theta} \sum_{(\mathbf{x}, \mathbf{y}) \in \mathcal{D}} \log \mathrm{Pr}_{\theta} (\mathbf{y}|\mathbf{x}),
\end{eqnarray}

\noindent 其中，$\mathrm{Pr}_{\theta}(\cdot)$ 表示由 $\theta$ 参数化的模型。

如果 $\mathbf{y}$ 只是最终结论，模型就会学会预测答案，而不显式生成相应推理路径。在基于推理的训练或调优中，改为将 $\mathbf{y}$ 定义为推理路径。如式 (\ref{eq:reasoning-path-notation}) 所示，可以将输出表示为由一系列推理步骤组成的路径 $\mathbf{y} = \{\bar{\mathbf{y}}_1,...,\bar{\mathbf{y}}_{n_p}\}$。于是，上述目标可重写为
\begin{eqnarray}
\hat{\theta} & = & \argmax_{\theta} \sum_{(\mathbf{x}, \{\bar{\mathbf{y}}_1,...,\bar{\mathbf{y}}_{n_p}\}) \in \mathcal{D}} \log \mathrm{Pr}_{\theta} (\bar{\mathbf{y}}_1,...,\bar{\mathbf{y}}_{n_p}|\mathbf{x}).
\end{eqnarray}

\noindent 这样，模型就会显式获得推理细节，并被训练为生成包含最终结论在内的完整推理路径。

在实践中，为将核心推理过程与输出的其他部分分开，推理模型往往采用特殊格式词元。一种广泛使用的做法，是将核心推理步骤包裹在 \texttt{<think>} 和 \texttt{</think>} 标签中。模型按这一格式生成文本时，标签内的推理步骤称为\mindex{思维链轨迹}（CoT trace），或简称\mindex{推理轨迹}（reasoning trace）。可以将思维链轨迹看作大语言模型的草稿纸。就像学生参加数学考试时，先在草稿纸上写中间计算，再将最终答案写到正式答卷上，大语言模型也使用 \texttt{<think>} 块来“展示解题过程”。

例如，采用这些标签组织的训练样例可以写成：

\vspace{0.5em}
\begin{tcolorbox}[frame empty]
\begingroup
\setlength{\leftskip}{2em}
\setlength{\rightskip}{2em}

$15\%$ 乘以 $80$ 是多少？

\vspace{0.3em}

\texttt{<think>}

为方便计算，可以将 $15\%$ 拆成 $10\%$ 和 $5\%$。

首先，$10\%$ 乘以 $80$ 是 $8$。

接着，$5\%$ 是 $10\%$ 的一半，即 $4$。

将两者相加：$8 + 4 = 12$。

\texttt{</think>}

答案是 12。

\endgroup
\end{tcolorbox}
\vspace{0.5em}

这里的问题是“$15\%$ 乘以 $80$ 是多少？”。输出包含五个步骤：步骤 1—4 构成思维链轨迹，步骤 5 为结论。\texttt{<think>} 和 \texttt{</think>} 都是有效词元，因此模型会学会生成它们，以标示核心推理步骤。在用户端，可以移除这些标签，呈现简洁的输出。需要注意，实际应用中的输出可能十分复杂，还包含与推理无关的内容。因此，引入 \texttt{<think>} 和 \texttt{</think>} 标签，有助于模型专注于求解推理问题。

推理数据也可以用多种其他格式表示。一种直接替代方式是自然语言的逐步格式：不使用特殊词元，而依靠“步骤 1：”“步骤 2：”这样的文本前缀，或先用“思考：”，后接“答案：”。此外，还可以通过按角色分离消息来组织推理数据。在现代聊天模板中，不再将推理与答案混在同一条消息内，而是将推理轨迹分配给内部角色，例如“thought”或“system\_scratchpad”，将最终结论分配给标准“assistant”角色。这样，就在系统层面分离了大语言模型的内部思考与面向用户的回答。

\subsection{过程监督}

如第 \ref{sec:reasoning-verification} 节所述，过程奖励模型是在推理阶段验证推理步骤的强大工具。一个自然延伸，是通过过程监督（process supervision），利用它指导大语言模型训练。

在标准结果监督（outcome supervision，如监督微调）中，只根据整条推理路径 $\mathbf{y}$ 的最终结果评估模型。如果模型生成了五个步骤，并在第三步犯下逻辑错误，基于结果的训练方法可能只是惩罚整个序列，模型必须自行猜测到底是哪一步导致失败。这就是\textbf{信用分配问题}（credit assignment problem）\cite{minsky:1961steps,sutton:1984temporal}。

过程监督通过在步骤层面施加训练信号来解决这一问题。模型不再只获得整条路径的单个奖励，而是对每个步骤 $\bar{\mathbf{y}}_t$ 都获得反馈。例如，设想教学生求解一个复杂代数方程。如果只看他写下的最终数字，然后说“错了”，学生并不知道，是理解错了公式、犯了简单加法错误，还是漏了一个负号。过程监督就像老师明确指向第三行，告诉学生：“你的逻辑在这里出了问题。”

过程监督通常表述为强化学习问题。将强化学习用于过程监督时，一般采用包含三个步骤的训练循环：

\begin{itemize}
\item \textbf{轨迹采样}：大语言模型接收输入 $\mathbf{x}$，生成完整推理路径 $\mathbf{y} = \{\bar{\mathbf{y}}_1, \dots, \bar{\mathbf{y}}_{n_p}\}$。这一过程可以重复多次，生成推理路径集合 $\mathcal{Y}$。
\item \textbf{验证}：奖励模型，即检查中间步骤的过程奖励模型，评估每条路径 $\mathbf{y}$，返回奖励分数 $R(\mathbf{x}, \mathbf{y})$。
\item \textbf{更新}：强化学习算法使用这些奖励更新模型参数 $\theta$。
\end{itemize}

一个简单训练目标，可以形式化为最大化期望奖励：
\begin{eqnarray}
J(\theta) & = & \mathbb{E}_{\mathbf{y} \sim \mathrm{Pr}_{\theta}(\cdot|\mathbf{x})} \left[ R(\mathbf{x}, \mathbf{y}) \right]. \label{eq:simple-rl-objective-for-reasoning}
\end{eqnarray}

\noindent 可以修改这一目标，使训练更加稳健，例如像近端策略优化（proximal policy optimization，PPO）那样，在目标中引入优势函数。

这看起来与第 \ChapterLLMTuning 章介绍的通用大语言模型对齐强化学习流程相同，但奖励的性质存在根本差别。在典型对齐任务中，奖励模型通常评估整段生成文本，只在末尾返回稀疏奖励。强化学习算法随后需要猜测，到底是哪些词元或片段对最终分数作出了贡献，因此会面临信用分配问题。

推理则是逐步展开的。因此，验证分数本质上就是步骤级信号，为简单起见，将其记为 $v_t = v(\mathbf{x}, \bar{\mathbf{y}}_{\le t})$。这一特征使推理任务十分适合现代强化学习。由于 $v_t$ 逐步生成，它自然地将稀疏奖励的语言生成任务转化为稠密奖励环境。

因此，主要技术问题是如何将这些自然产生的步骤级分数 $v_t$，恰当地转化为算法要求的正式训练信号，例如结构化的步骤奖励 $r_t$ 或优势 $A_t$。根据过程奖励模型的输出格式和强化学习框架，下面考虑三种方法。

\subsubsection{使用二元验证分数训练}

某些情况下，过程奖励模型充当严格规则检查器，例如数学计算器或代码编译器。它只返回二元分数：$v_t \in \{0, 1\}$。分数为 $1$，表示给定前面步骤时，第 $t$ 步的逻辑正确；$0$ 则表示该步骤出现逻辑错误。

处理二元分数时，思路十分直接：只要一个步骤错误，整条推理路径就无效。如果希望为整条路径分配序列级奖励 $R(\mathbf{x}, \mathbf{y})$，只需将所有步骤分数相乘：
\begin{eqnarray}
R(\mathbf{x}, \mathbf{y}) & = & \prod_{t=1}^{n_p} v_t.
\end{eqnarray}

\noindent 因此，只有当每一步的验证分数都是 $1$ 时，模型才能获得正奖励。

另一种方式是，如果强化学习算法支持步骤级奖励，记为 $r_t$，就可以在失败发生之处准确惩罚模型。例如，给予奖励 $1$，条件是步骤保持正确，即 $v_t = 1$；模型一旦犯错，即 $v_t = 0$，就赋予负奖励并停止生成。


\subsubsection{使用实值验证分数训练}

在多数神经过程奖励模型中，验证器输出连续实数 $v_t \in \mathbb{R}$。它可以是从 $0$ 到 $1$ 的概率，表示推理路径正确的置信度，也可以是未归一化的标量值。

使用连续验证分数时，最常见的技术是\textbf{奖励塑形}（reward shaping）。不直接使用绝对分数，而将即时步骤奖励 $r_t$ 定义为当前步骤与前一步骤之间的边际差值：
\begin{eqnarray}
r_t & = & v_t - v_{t-1}.
\end{eqnarray}

这一方法很直观。设想一个过程奖励模型正在估计最终答案正确的概率。步骤 1 的置信度为 $0.4$，步骤 2 包含精彩的演绎，将置信度提高到 $0.9$，则步骤 2 的奖励是 $r_2 = 0.9 - 0.4 = +0.5$。模型因这一具体逻辑推进获得奖励。如果步骤 3 引入计算错误，使置信度下降到 $0.1$，奖励就是 $r_3 = 0.1 - 0.9 = -0.8$。

对整个序列中的塑形奖励求和，会得到一个“裂项相消求和”（telescoping sum），构成序列级奖励 $R(\mathbf{x}, \mathbf{y})$：
\begin{eqnarray}
R(\mathbf{x}, \mathbf{y}) & = & \sum_{t=1}^{n_p} r_t \nonumber \\
& = & (v_1 - v_0) + (v_2 - v_1) + \dots + (v_{n_p} - v_{n_p-1}) \nonumber \\
& = & v_{n_p} - v_0.
\end{eqnarray}

\noindent 这里通常假定 $v_0$ 为 0。在这一推导中，中间分数相互抵消，意味着整个序列的奖励与最后一步的评估保持一致。

要理解这种逐步稠密奖励如何真正训练模型，需要将其与强化学习目标联系起来。如式 (\ref{eq:simple-rl-objective-for-reasoning}) 所示，目标是最大化所有可能生成路径上的期望总奖励。为提高训练效率并恰当分配信用，通常采用因果性原则：第 $t$ 步的动作只能影响该步骤之后获得的奖励。因此，不使用整个序列的奖励评价某一步，而是使用未来回报 $R_t$，定义为从第 $t$ 步到序列末尾的奖励之和：
\begin{eqnarray}
R_t & = & \sum_{k=t}^{n_p} r_k \nonumber \\
& = & v_{n_p} - v_{t-1}.
\end{eqnarray}

最后，可以使用替代损失函数训练模型：
\begin{eqnarray}
\mathcal{L}(\theta) & = & - \sum_{t=1}^{n_p} \log \mathrm{Pr}_{\theta}(\bar{\mathbf{y}}_t | \mathbf{x}, \bar{\mathbf{y}}_{<t}) \cdot R_t.
\end{eqnarray}

这一损失函数提供步骤级反馈机制。对于任何生成的推理路径，模型都会考察未来回报 $R_t$，并在每一步 $t$ 利用这一回报。如果 $R_t$ 为正，损失函数会促使模型提高生成该步骤的概率 $\log \mathrm{Pr}_\theta(\cdot)$；如果 $R_t$ 为负，它就作为惩罚，促使模型今后避免生成这一步骤。


\subsubsection{将验证纳入优势函数}

许多强化学习算法，例如 PPO，使用优势函数 $A_t$ 优化模型，衡量一个动作比预期好多少。可以不把验证分数直接当作原始奖励，而是将其纳入优势计算。

一种方法是将步骤级验证引入时序差分误差。如果过程奖励模型被训练为估计期望未来奖励，就可以将其视为标准演员—评论家强化学习中的价值基线。使用时序差分误差计算优势 $A_t$：
\begin{eqnarray}
A_t & = & \tilde{r}_{t} + \gamma v_t - v_{t-1},
\end{eqnarray}

\noindent 其中，$\tilde{r}_{t}$ 是环境提供的任意外部奖励，中间思考步骤中通常为 $0$，$\gamma$ 为折扣因子。

另一种方法是改进组相对优势的设计。在\mindex{组相对策略优化}（group relative policy optimization，\mindex{GRPO}）中，不再使用绝对验证分数 \cite{shao-etal:2024deepseekmath}，而是通过比较一组推理路径计算优势。对于输入 $\mathbf{x}$，大语言模型生成一组不同路径。验证器给出分数 $v_t^{(i)}$，表示第 $t$ 个步骤在第 $i$ 条路径中的得分。优势 $A_t^{(i)}$ 通过累加从第 $t$ 步到路径末尾的标准化分数得到：
\begin{eqnarray}
A_t^{(i)} & = & \sum_{k = t}^{n_p^{(i)}} \tilde{v}_k^{(i)},
\end{eqnarray}

\noindent 其中，$n_p^{(i)}$ 是第 $i$ 条路径的总步骤数，$\tilde{v}_k^{(i)}$ 为标准化分数，定义为：
\begin{eqnarray}
\tilde{v}_k^{(i)} & = & \frac{v_k^{(i)} - \mu}{\sigma},
\end{eqnarray}

\noindent 其中，$\mu$ 和 $\sigma$ 是组内所有路径全部步骤分数的均值和标准差。例如，模型探索了求解一个方程的 5 种不同方式，如果路径 \#3 执行的某一步，其分数显著高于这些路径中所有步骤的平均值，那么路径 \#3 就会获得正优势。这里，组均值 $\mu$ 充当基线，无需再训练和运行独立的价值网络。


无论使用哪一种方法，计算出优势 $A_t$ 或 $A_t^{(i)}$ 后，都可以将其作为损失函数中的缩放权重。基于优势的损失函数，其简化形式可写为\footnote{PPO、GRPO 等进阶算法通常会在目标中加入裁剪机制，限制新旧策略之间的概率比，避免模型参数在一次更新中变化过大，但核心数学原理相同。}：
\begin{eqnarray}
\mathcal{L}(\theta) & = & - \sum_{t=1}^{n_p} \log \mathrm{Pr}_{\theta}(\bar{\mathbf{y}}_t | \mathbf{x}, \bar{\mathbf{y}}_{<t}) \cdot A_t.
\end{eqnarray}

使用优势 $A_t$ 而非原始回报 $R_t$ 的好处，是降低方差。通过减去基线，无论是估计价值 $v_{t-1}$ 还是组均值 $\mu$，优势函数都提供相对表现的度量。它不仅告诉模型某个推理步骤在逻辑上是否合理，还告诉模型这一表现是否好于预期。这种相对反馈，能够避免过于简单或过于困难的问题主导模型训练。

\subsection{用于推理的同策略蒸馏}

过程监督虽然提供了逐步奖励信号来引导推理，但仍依靠一个标量值告诉模型表现如何。为了提供更丰富的训练信号，可以采用\mindex{同策略蒸馏}（on-policy distillation，\mindex{OPD}）。这一方法通常用于大语言模型后训练，与步骤级监督的目标十分契合：让训练信号尽可能稠密。

标准知识蒸馏可以视为一种\mindex{异策略蒸馏}（off-policy distillation）。较小的“学生”模型通过模仿更强“教师”模型生成的推理路径学习，训练数据集由预先记录的固定路径构成。然而，推理过程漫长且脆弱。如果实际生成时，学生犯了小错，或选择了略有不同但仍有效的方法，就会进入训练数据中从未见过的状态。由于学生不知道如何从错误中恢复，推理过程可能崩溃。这一现象称为暴露偏差（exposure bias）或分布偏移（distribution shift）。

同策略蒸馏的核心思想，是让学生模型从自己生成的轨迹中学习，而不是从教师模型采样的轨迹中学习。形式化地，令教师模型为 $\mathrm{Pr}^{t}(\cdot)$，学生模型以 $\theta$ 为参数，记为 $\mathrm{Pr}_{\theta}^{s}(\cdot)$。同策略蒸馏的典型过程包含两步：

\begin{itemize}
\item \vspace{0.5em} \textbf{学生轨迹采样}：给定输入 $\mathbf{x}$，学生模型使用当前参数 $\theta$，生成推理路径 $\mathbf{y} = \{\bar{\mathbf{y}}_1, \dots, \bar{\mathbf{y}}_{n_p}\}$。由于路径由学生自己生成，它会自然探索自身的推理行为，包括轻微偏离、采用替代方法，甚至犯下教师绝不会犯的错误。
\item \vspace{0.3em} \textbf{教师指导}：随后，将学生生成的路径输入教师模型。在每一步 $t$，教师模型不提供标量奖励，而是计算下一步的概率分布 $\mathrm{Pr}^{t}(\cdot | \mathbf{x}, \bar{\mathbf{y}}_{<t})$。
\end{itemize}
\vspace{0.5em}

理想情况下，更新学生模型，使其步骤级概率分布与教师模型生成的分布对齐。目标是在学生模型生成的推理路径上，最小化两个分布之间的期望 KL 散度：
\begin{eqnarray}\mathcal{L}_{\mathrm{step}}(\theta) & = & \mathbb{E}_{\mathbf{y} \sim \mathrm{Pr}_{\theta}^{s}(\cdot|\mathbf{x})} \left[ \sum_{t=1}^{n_p} D_{\mathrm{KL}} \left( \mathrm{Pr}_{\theta}^{s}(\cdot | \mathbf{x}, \bar{\mathbf{y}}_{<t}) \,\|\, \mathrm{Pr}^{t}(\cdot | \mathbf{x}, \bar{\mathbf{y}}_{<t}) \right) \right].
\end{eqnarray}

\noindent 不过，在实践中计算这一步骤级分布并不可行。由于推理步骤 $\bar{\mathbf{y}}_t$ 通常包含多个词元，所有可能多词元步骤上的概率分布处于指数级庞大的空间中。

为解决这一问题，退回到词元层面运行的标准同策略蒸馏方法。将生成的推理路径视为词元序列 $\mathbf{y} = \{y_1,\dots, y_n\}$，重新表述目标，在词表上逐个词元计算 KL 散度：
\begin{eqnarray}
\mathcal{L}_{\mathrm{token}}(\theta) & = & \mathbb{E}_{\mathbf{y} \sim \mathrm{Pr}_{\theta}^{s}(\cdot|\mathbf{x})} \left[ \sum_{i=1}^{n} D_{\mathrm{KL}} \left( \mathrm{Pr}_{\theta}^{s}(\cdot | \mathbf{x}, \mathbf{y}_{<i}) \,\|\, \mathrm{Pr}^{t}(\cdot | \mathbf{x}, \mathbf{y}_{<i}) \right) \right].
\end{eqnarray}

\noindent 随后，通过采样近似这一目标。给定输入 $\mathbf{x}$，只需对学生模型执行轨迹采样，生成一条或一小批推理路径，再在这些采样轨迹上计算词元级 KL 散度。

图 \ref{fig:comparision-of-sft-ps-opd} 比较了同策略蒸馏、监督微调和过程监督。为了进一步理解几种方法的区别，可以设想一个学生正在解一道复杂几何题：

\begin{figure}[!t]
\centering
% !Mode:: "TeX:UTF-8"
% !TEX encoding = UTF-8 Unicode

\begin{tikzpicture}
	
\def\ssep{1.0cm}

\tikzstyle{mnode} = [minimum width=1.8cm,minimum height=0.8cm,draw,thick,inner sep=3pt]
\tikzstyle{wnode} = [minimum width=0.4cm,minimum height=0.4cm,inner sep=2pt]

%%% sft
\begin{scope}

\node [mnode,anchor=west] (student) at (0,0) {\small{学生模型}};
\node [mnode,anchor=south,align=left] (teacher) at ([yshift=1.5cm]student.north) {\small{人工/}\\[-3pt]\small{教师模型}};

\node [anchor=west,wnode] (w1) at ([xshift=3cm]teacher.east) {\small{This}};
\node [anchor=center,wnode] (w2) at ([xshift=\ssep]w1.center) {\small{cat}};
\node [anchor=center,wnode] (w3) at ([xshift=\ssep]w2.center) {\small{is}};
\node [anchor=center,wnode] (w4) at ([xshift=\ssep]w3.center) {\small{cute}};
\node [anchor=center,wnode] (w5) at ([xshift=\ssep]w4.center) {\small{.}};

\begin{pgfonlayer}{background}
\node [draw, rounded corners=1mm]  [fit = (w1) (w5)] (ty) {};
\end{pgfonlayer}

\draw [->] ([xshift=1pt]teacher.east) -- ([xshift=-1pt]ty.west) node [pos=0.5,above] {\footnotesize{标注/生成}};
\node [anchor=south] (tylabel) at (ty.north) {\footnotesize{目标输出 ($\mathbf{y}^*$)}};

\draw [<->,blue,thick,dashed] ([xshift=-1cm,yshift=-1pt]ty.south) .. controls +(south:1.7cm) and +(east:2cm) .. ([xshift=1pt]student.east) node [pos=0.5,right,align=left,xshift=0.8cm,yshift=-0.1cm] {\footnotesize{教师强制}\\[-1pt] \footnotesize{$\max\limits_{\theta} \Pr_{\theta}^s(\mathbf{y}^*|\mathbf{x})$}};

\node [anchor=north] (caption) at ([xshift=3cm,yshift=-0.5cm]student.south east) {\small{(a) 监督微调}};

\end{scope}

%%% process supervision
\begin{scope}[yshift=-5.2cm]

\node [mnode,anchor=west] (student) at (0,0) {\small{学生模型}};

\node [anchor=west,wnode] (w1) at ([xshift=3cm]student.east) {\small{The}};
\node [anchor=center,wnode] (w2) at ([xshift=\ssep]w1.center) {\small{cat}};
\node [anchor=center,wnode] (w3) at ([xshift=\ssep]w2.center) {\small{is}};
\node [anchor=center,wnode,fill=gray!30] (w4) at ([xshift=\ssep,yshift=-0.05cm]w3.center) {\small{lovely}};
\node [anchor=center,wnode] (w5) at ([xshift=\ssep,yshift=0.05cm]w4.center) {\small{.}};

\begin{pgfonlayer}{background}
\node [draw, rounded corners=1mm]  [fit = (w1) (w5)] (sy) {};
\end{pgfonlayer}

\draw [->] ([xshift=1pt]student.east) -- ([xshift=-1pt]sy.west) node [pos=0.5,above] {\footnotesize{生成}};
\node [anchor=north] (sylabel) at (sy.south) {\footnotesize{学生模型输出 ($\mathbf{y}$)}};

\node [mnode,anchor=south] (verifier) at ([yshift=2.0cm]w4.north) {\small{验证器}};

\draw [->] ([yshift=-1pt]verifier.south) -- ([yshift=0pt]w4.north) node (verificationlabel) [pos=0.4,left,align=right] {\footnotesize{生成验证}\\[-2pt]\footnotesize{分数 $v_t$，并计算}\\[-2pt]\footnotesize{优势 $A_t$}};

\draw [<->,blue,thick,dashed] ([yshift=1pt]student.north) .. controls +(north:1cm) and +(west:2cm) .. ([yshift=0.3cm]verificationlabel.west) node [pos=0.3,above,yshift=0.3cm,align=right] {\footnotesize{强化学习训练}\\[-1pt]\footnotesize{$\max\limits_{\theta} \sum_{t=1}^{n_p} \log \mathrm{Pr}_{\theta}(\bar{\mathbf{y}}_t | \mathbf{x}, \bar{\mathbf{y}}_{<t}) \cdot A_t$}};

\node [anchor=north] (caption) at ([xshift=3cm,yshift=-0.5cm]student.south east) {\small{(b) 过程监督}};

\end{scope}

%%% odp
\begin{scope}[yshift=-10.4cm]

\node [mnode,anchor=west] (student) at (0,0) {\small{学生模型}};

\node [anchor=west,wnode] (w1) at ([xshift=3cm]student.east) {\small{The}};
\node [anchor=center,wnode] (w2) at ([xshift=\ssep]w1.center) {\small{cat}};
\node [anchor=center,wnode] (w3) at ([xshift=\ssep]w2.center) {\small{is}};
\node [anchor=center,wnode,fill=gray!30] (w4) at ([xshift=\ssep,yshift=-0.05cm]w3.center) {\small{lovely}};
\node [anchor=center,wnode] (w5) at ([xshift=\ssep,yshift=0.05cm]w4.center) {\small{.}};

\begin{pgfonlayer}{background}
\node [draw, rounded corners=1mm]  [fit = (w1) (w5)] (sy) {};
\end{pgfonlayer}

\draw [->] ([xshift=1pt]student.east) -- ([xshift=-1pt]sy.west) node [pos=0.5,above] {\footnotesize{生成}};
\node [anchor=north] (sylabel) at (sy.south) {\footnotesize{学生模型输出 ($\mathbf{y}$)}};

\node [mnode,anchor=south] (teacher) at ([yshift=2.0cm]w4.north) {\small{教师模型}};

\node [anchor=south,minimum width=0.2cm,minimum height=0.4cm,fill=orange!30] (sh3) at ([yshift=0.5cm]w4.north) {};
\node [anchor=south east,minimum width=0.2cm,minimum height=0.3cm,fill=orange!30] (sh2) at ([xshift=-1pt]sh3.south west) {};
\node [anchor=south east,minimum width=0.2cm,minimum height=0.2cm,fill=orange!30] (sh1) at ([xshift=-1pt]sh2.south west) {};
\node [anchor=south west,minimum width=0.2cm,minimum height=0.2cm,fill=orange!30] (sh4) at ([xshift=1pt]sh3.south east) {};
\node [anchor=south west,minimum width=0.2cm,minimum height=0.3cm,fill=orange!30] (sh5) at ([xshift=1pt]sh4.south east) {};

\node [anchor=south,minimum width=0.2cm,minimum height=0.3cm,fill=purple!30] (th3) at ([yshift=1.2cm]w4.north) {};
\node [anchor=south east,minimum width=0.2cm,minimum height=0.2cm,fill=purple!30] (th2) at ([xshift=-1pt]th3.south west) {};
\node [anchor=south east,minimum width=0.2cm,minimum height=0.4cm,fill=purple!30] (th1) at ([xshift=-1pt]th2.south west) {};
\node [anchor=south west,minimum width=0.2cm,minimum height=0.3cm,fill=purple!30] (th4) at ([xshift=1pt]th3.south east) {};
\node [anchor=south west,minimum width=0.2cm,minimum height=0.2cm,fill=purple!30] (th5) at ([xshift=1pt]th4.south east) {};

\draw [->] ([yshift=1pt]w4.north) -- ([yshift=-1pt]sh3.south);
\draw [->] ([yshift=-1pt]teacher.south) -- ([yshift=-0.3cm]teacher.south);

\node [anchor=south west,align=left,inner sep=2pt] (tdislabel) at ([yshift=-0.1cm,xshift=0.1cm]th5.south east) {\tiny{教师预测的}\\[-6pt]\tiny{分布 $\mathrm{Pr}^{t}(\cdot | \mathbf{x}, \mathbf{y}_{<i})$}};
\node [anchor=south west,align=left,inner sep=2pt] (sdislabel) at ([yshift=-0.1cm,xshift=0.1cm]sh5.south east) {\tiny{学生预测的}\\[-6pt]\tiny{分布 $\mathrm{Pr}^{s}(\cdot | \mathbf{x}, \mathbf{y}_{<i})$}};

\draw [<->,blue,thick,dashed] ([yshift=1pt]student.north) .. controls +(north:1.2cm) and +(west:3cm) .. ([xshift=-1pt]th1.west) node [pos=0.2,above,yshift=0.4cm,xshift=1.2cm,align=right] {\footnotesize{最小化 KL 散度}\\[-1pt]\footnotesize{$\min\limits_{\theta} \sum_{i=1}^{n} D_{\mathrm{KL}} \left( \mathrm{Pr}_{\theta}^{s}(\cdot | \mathbf{x}, \mathbf{y}_{<i}) \parallel \mathrm{Pr}^{t}(\cdot | \mathbf{x}, \mathbf{y}_{<i}) \right)$}};


\node [anchor=north] (caption) at ([xshift=3cm,yshift=-0.5cm]student.south east) {\small{(c) 同策略蒸馏}};

\end{scope}
	
	
\end{tikzpicture}

\caption{监督微调、过程监督与同策略蒸馏的比较。在监督微调 (a) 中，学生模型采用异策略训练，最大化人类或教师模型提供的标准目标输出的似然。在过程监督 (b) 中，学生模型生成自己的推理路径，由验证器评估这些自行生成的步骤，并计算优势。在同策略蒸馏 (c) 中，学生同样生成自己的路径，但教师模型在每一步提供稠密的词元级概率分布，而不是标量奖励。通过最小化学生预测分布与教师分布之间的 KL 散度，更新学生模型。}
\label{fig:comparision-of-sft-ps-opd}
\end{figure}

\begin{itemize}
\item \vspace{0.5em} \textbf{监督微调}：就像给学生一本教科书，里面都是老师写好的正确解答，让学生记忆。如果学生决定画一条与书中不同的辅助线，教科书就无法指导他如何从这个新起点继续求解。
\item \vspace{0.5em} \textbf{过程监督}（使用过程奖励模型的强化学习）：就像一位严格的阅卷老师，逐步检查学生的解题过程，并说“步骤 1 加 1 分”“步骤 2 减 1 分”。它告诉学生哪里出了错，但不会告诉他如何修正，或下一步应该怎样正确进行。
\item \textbf{同策略蒸馏}：就像一位专业辅导老师坐在学生身旁。学生画出那条不常见的辅助线，进入标准教科书中没有的状态时，老师会看着学生当前的解题进度，说：“好，既然你把线画在这里，接下来最好计算这个特定角度。”
\end{itemize}

同策略蒸馏的监督形式，是词表上的完整概率分布，因此学生获得的学习信号十分稠密、丰富。它不仅学到单个最佳后续步骤，也学到其他有效备选方案的概率。这种丰富反馈非常适合推理，因为它教会学生模型如何从偏离中恢复，并在整个输出序列中保持逻辑一致，而不是把推理依据当作孤立部分。不过，同策略蒸馏的计算成本很高：训练时需要在运行学生模型的同时，运行大型教师模型，实时计算概率分布。这比标准异策略蒸馏更昂贵；后者只需生成一次教师模型输出，再将其存入静态数据集。

\subsection{扩展推理训练规模}

扩展是大语言模型时代的核心策略。面对复杂推理任务，扩展训练过程是一个自然思路。这里简要讨论推理大语言模型训练扩展的几个关键方面：大规模强化学习、合成数据生成，以及长上下文扩展。

\subsubsection{大规模强化学习}

强化学习已成为提升大语言模型“系统 2”推理能力的主要方式。例如，许多强大推理模型的成功，来自强化学习规模的扩展，而不只是监督微调。在基于大规模强化学习的训练中，观察到一种有趣现象，即“顿悟时刻”（aha moment）的出现~\cite{deepseek:2025r1}。无需显式人工标注，模型就能学会停下来，意识到自己犯错，甚至回溯并尝试另一种方法。仅依赖有限标注数据和标准监督训练，很难激发这些能力。随着训练计算量扩展，模型可以为更难的问题分配更长的思维链轨迹，从而自行发现如何深入思考~\cite{snell-etal:2024scaling}。

近期大规模实证研究表明，强化学习计算量与推理表现之间的关系，可以用 S 形曲线建模 \cite{khatri-etal:2025art,tan-etal:2026scaling}。这一缩放定律的关键发现是，不同训练选择对曲线的影响不同。例如，提高大语言模型的基础能力，如采用更大架构或混合专家（mixture of experts，MoE），以及扩展上下文窗口长度，可以提高性能上限。相比之下，改变损失函数、调整优势归一化等算法改进，主要能更好地利用现有计算预算，却不会显著改善最终性能。

扩展强化学习十分困难。首先，计算成本很高。对于每个训练样本，系统通常都要执行轨迹采样，即运行模型进行推理，然后额外调用奖励模型或另一个大语言模型评估生成输出。随着训练样本增加，这一成本会显著高于标准监督微调。计算负担也使实验和试错变得更困难。此外，强化学习方法通常复杂，需要极为细致的调优才能取得良好结果。例如，它们对超参数设置和训练数据很敏感，因此相对脆弱。在大规模训练中，确保系统稳定运行，例如让训练损失稳定下降，既是机器学习问题，也是工程挑战。

\subsubsection{生成更多推理合成数据}

复杂多步推理的高质量人工标注稀缺，因此扩大训练数据规模往往依赖合成数据生成。研究者不再手动编写逐步证明，而是用强大的大语言模型生成多样化推理轨迹，再经过筛选和整理，用来训练新模型~\cite{zelikman-etal:2022star,mukherjee-etal:2023orca}。

为生成推理合成数据，一种常见方式是针对给定种子问题，从教师模型采样多条推理路径~\cite{wang-etal:2023self}。然后，使用验证器或执行环境评估这些路径，只保留逻辑步骤合理、最终答案正确的路径。此外，还可以用模型生成已有问题的变体，例如修改数学或编程任务中的数值和结构约束，拓展训练数据的分布。

\subsubsection{扩展上下文长度}

随着推理模型处理的问题越来越复杂，推理路径长度已成为主要瓶颈。推理对大上下文窗口的需求有几个来源。第一，为激发深入的系统 2 思维，输入提示本身往往就需要大量上下文，包括详细逻辑约束、复杂环境规则，或多个较长的少样本思维链示例。第二，复杂任务的生成推理过程，在得出最终结论之前，很容易延伸到数千、甚至数万词元。第三，大语言模型智能体的近期进展，也增加了对推理能力的需求。智能体推理（agentic reasoning）往往跨越多轮交互，进一步加大了对长上下文窗口的需求。

这种需求直接推动大语言模型行业开发支持超大上下文窗口的模型。例如，许多模型已支持最多 128K 词元的上下文，更先进的模型，如 Anthropic 的 Claude 3.5~\cite{anthropic:2024claude35} 和 Google 的 Gemini 1.5~\cite{gemini:2024gemini}，支持数十万、甚至一百万词元。关于长上下文大语言模型的详细讨论，感兴趣的读者可以参阅第 \ChapterLLM\ 章。

不过，为推理处理长上下文，带来了超出一般上下文压缩的新挑战。在长篇小说摘要、文档要点提取等任务中，模型可以压缩信息，损失少量细节；但推理需要严格、高精度的逻辑连续性。每个中间步骤，都是相互高度依赖的链条中不可或缺的一环。如果模型“忘记”提示中定义的特定假设，或在执行步骤 50 时，无法追踪步骤 2 算出的变量，整个推导就会崩溃。因此，为推理扩展上下文，不只是增大键值缓存（key-value cache，KV 缓存）以存储更多词元，还要求模型能够可靠地跨长距离检索信息，并保持严格的逻辑一致性。

%----------------------------------------------------------------------------------------
% NEW SECTION
%----------------------------------------------------------------------------------------

\section{评估推理大语言模型}

本节讨论如何评估推理大语言模型。这与评估标准大语言模型相似 \cite{chang-etal:2024survey}，但推理问题十分复杂，往往需要更详细的评估和其他特别考虑。

\subsection{基于结果与基于过程的评估}

评估大语言模型的推理表现时，最简单、最常用的方法是\textbf{基于结果的评估}（outcome-based evaluation），只检查模型生成的最终答案。如果最终答案与标准答案一致，就认为模型作答正确。GSM8K \cite{cobbe-etal:2021training}、MATH \cite{hendrycks-etal:2measuring} 等常用推理基准，都广泛采用这一方法。不过，基于结果的评估存在“答案正确、理由错误”的问题，也称为假阳性。模型可能生成有缺陷的逻辑步骤，却碰巧得出正确的最终数值。如果只检查最终结果，就可能把运气好误认为推理能力强。

为了更详细地评估，可以采用\textbf{基于过程的评估}（process-based evaluation）：不只检查最终答案，而是逐步检查整条推理路径 \cite{wang-etal:2024math,xia-etal:2025evaluating}。只有中间逻辑与最终结论都合理，才认为模型作答正确。这是一种理想方法，但实践中很难实施。主要挑战在于，步骤级“标准数据”稀缺，构建成本高。此外，任何复杂问题都可能存在指数级数量的有效推理路径，无法为模型可能采用的每条路径，人工列出所有正确中间步骤。

为了使基于过程的评估更实用，可以缩小范围，只评估局部推理决策。一种方式是将问题简化为下一步评估任务：给定一组正确的前序步骤，检查推理模型能否正确选择或生成下一逻辑步骤。为此，可以直接使用过程奖励模型，评估模型生成的每一步质量。这也与奖励模型评估有关 \cite{lambert-etal:2025rewardbench}：不考察推理模型输出，而是评估过程奖励模型的质量。此外，还可以考察推理模型的基础能力，例如识别给定推理路径中错误的能力。这也可以视为对元推理能力的检验：不只要求模型进行推理，还要求它评估推理过程本身 \cite{zeng-etal:2024mr}。

实践中，基于过程的评估并不依赖单一标准。可以从逻辑一致性、连贯性和客观性等不同角度，判断推理步骤的质量 \cite{lee-and-hockenmaier:2025evaluating}。不过，这种评估虽然详细得多，标注成本也高得多，而且通常不够稳健。因此，大多数实际研究仍依赖更简单的基于结果的评估。

\subsection{评估指标}

下面列举一些用于评估推理大语言模型的指标。

\begin{itemize}
\item \vspace{0.5em} \textbf{精确匹配与准确率：}对于具有明确最终答案的任务，例如数学应用题，这是最简单、使用最广泛的指标。假设数据集 $\mathcal{D}$ 包含 $N$ 个问题。对于第 $i$ 个问题，从生成的推理路径中提取最终答案 $\hat{a}_i$，与标准答案 $a_i^*$ 比较。准确率定义为：
\begin{eqnarray}
\text{准确率} & = & \frac{1}{N} \sum_{i=1}^N \mathbbm{1}(\hat{a}_i = a^*_i).
\end{eqnarray}

\noindent 实践中，精确匹配，即检查字符串是否完全相同，可能过于严格。例如，即使“$0.5$”与标准答案“$1/2$”相等，仍会被拒绝。因此，研究者通常在比较之前，先用基于规则的脚本，将答案归一化为规范形式。

\item \vspace{0.3em} \textbf{Pass@$k$：}大语言模型具有概率性，一次解码可能因细小错误而推理失败。Pass@$k$ 衡量允许模型生成 $k$ 条独立推理路径时，至少找到一个正确答案的概率 \cite{chen-etal:2021evaluating}。这一指标尤其常用于代码生成和复杂数学任务。对于第 $i$ 个问题，若它属于数据集 $\mathcal{D}$，采样 $N_s \ge k$ 条路径，其中 $C_i$ 条包含正确答案，则 Pass@$k$ 分数定义为：
\begin{eqnarray}
\text{Pass}@k & = &  \frac{1}{N} \sum_{i=1}^N \left( 1 - \frac{\binom{N_s-C_i}{k}}{\binom{N_s}{k}} \right).
\end{eqnarray}

\noindent 要生成多条不同推理路径，可以使用第 \ref{sec:sampling-reasoning-paths} 节介绍的采样方法。

\item \vspace{0.3em} \textbf{Cons@$k$（Consensus@$k$）：}Pass@$k$ 需要预言机判断 $k$ 个采样答案各自是否正确，Cons@$k$ 则不识别单个样本的正确性，而是构造一个共识答案。具体而言，生成 $k$ 条推理路径，提取最终答案，再通过多数投票选出最常见的答案 \cite{wang-etal:2023self}。共识答案记为 $\hat{a}_i^{\text{maj}}$，对于第 $i$ 个问题，其定义为：
\begin{eqnarray}
\hat{a}_i^{\text{maj}} & = & \argmax_{a_i} \sum_{j=1}^k \mathbbm{1}(\hat{a}_i^j = a_i),
\end{eqnarray}

\noindent 其中，$\hat{a}_i^j$ 是从第 $j$ 条推理路径中提取的最终答案。然后，将 $\hat{a}_{\text{maj}}$ 与标准答案比较，得到最终准确率：
\begin{eqnarray}
\text{Cons}@k & = & \frac{1}{N} \sum_{i=1}^N \mathbbm{1}(\hat{a}_i^{\text{maj}} = a^*_i).
\end{eqnarray}

Cons@$k$ 为实际推理系统提供了一种实用评估方式，衡量模型能否稳定地得到正确答案。Pass@$k$ 评估是否至少有一个采样回答正确，Cons@$k$ 则偏好在多条采样推理路径中持续获得支持的答案，强调推理的可靠性。

\item \vspace{0.3em} \textbf{以大语言模型为评判者（LLM-as-a-Judge）：}对于智能体规划、复杂开放域问题等开放式推理任务，基于规则的精确匹配往往失效。这时，可以使用大语言模型充当自动评判者。例如，可以提示评判模型评估生成的中间步骤，按预定义标准打分。还可以用它比较两个推理模型：向评判者展示两个模型对同一问题的输出，由其决定哪条推理路径更好，或是否平局。胜率则反映某个模型被偏好的次数比例。
\end{itemize}
\vspace{0.5em}

\subsection{尚待解决的问题}

一个问题是数据污染。语言模型在越来越庞大、从互联网收集的数据集上训练，要将测试数据排除在训练过程之外，变得十分困难。如果模型在预训练中已经见过 MATH、GSM8K 等推理基准的问题，就可能只是在复述记住的解答，而非即时推理。这使人们难以判断，较高基准分数反映的是真正的问题求解能力，还是良好的记忆。为此，研究社区应探索问题持续更新的新基准，以及从不公开发布的私有留出集。

另一个问题是稳健性。当前推理模型有时很脆弱：模型可能完美解出一道复杂数学题，但同一道题只要略微改写，或在提示中加入无关细节，就会完全失败。这种敏感性说明，模型有时依赖表层文本模式，并未真正理解背后的逻辑。要实现更稳健的推理，需要大语言模型能够轻松适应新任务，乃至更有挑战性的任务。

此外，大语言模型还可能生成存在幻觉或不忠实的推理路径。例如，它们可能虚构数学规则来得出结论。有时，生成文本看起来像完美的逻辑演绎，但模型实际上是通过隐藏机制得到答案，例如幸运猜中或使用记住的捷径。这种不忠实性表明，模型输出文本与实际内部表示之间存在差距。在强化学习设定中，这往往与奖励投机有关。如果只对最终答案给予奖励，模型常会学会用有缺陷的逻辑欺骗奖励系统。因此，在强化学习中引入过程监督等逐步监督，检查各中间步骤是否真正正确，很有价值。

第四个问题是基准饱和。推理大语言模型进步迅速，标准评估基准很快就会过时。仅在一两年前还被视为极其困难的测试，模型现在已经接近满分。这种饱和使人们难以比较新模型，或衡量领域的真正进展。因此，需要开发更困难的测试，评估模型在全新问题上进行真正逻辑推导的能力。人类水平的推理往往是开放且交互式的，因此另一种测试方式，可能是将推理模型用于自主智能体，让其探索环境、使用外部工具，并从自身错误中恢复。

%----------------------------------------------------------------------------------------
% NEW SECTION
%----------------------------------------------------------------------------------------

\section{补充讨论}

下面讨论与推理大语言模型有关的一些有趣问题。

\subsection{外部推理与内部推理}

前面的讨论假定，推理过程可以由思维链等文本推理路径很好地刻画。这引出了一个根本问题：生成的思维链文本，是否真的反映了模型如何推理？为此，需要区分外显推理与真正的内部“思考”过程。

提示大语言模型逐步思考，就是迫使它将计算外显为离散词元。然而，不能简单地把这一文本序列视为实际推理过程。一方面，模型可能生成详细、逻辑合理的思维链，但真正决策可能早已隐式发生在最初几个词元的隐表示中。另一方面，模型可能生成错误或毫无意义的思维链，却仍输出正确的最终答案。这说明，内部表示可能以不同于生成文本所暗示的方式追踪任务。

研究者已经表明，思维链解释可能具有误导性 \cite{turpin-etal:2023language,chen-etal:2025reasoning}。例如，尽管提示中已经提到正确答案，大语言模型仍可能倾向于生成一条逻辑思维链，导向错误答案。为揭示模型实际如何思考，研究者开始开展神经元激活分析和机制可解释性（mechanistic interpretability）研究，目标是理解生成过程中隐层内部发生了什么。例如，在 \citet{frasertaliente-etal:2026natural} 的研究中，使用自然语言自编码器，将激活转换成自然语言文本。自然语言自编码器包含激活语言化器和激活重构器，能够将内部残差流激活直接映射为人类可读文本，再映射回去。这样，就能通过解码这些激活，读取大语言模型的思考。

另一类研究关注不使用文本的推理，称为\mindex{潜空间推理}（latent reasoning）或\mindex{隐式推理}（implicit reasoning）。如前面定义，标准推理路径由一系列文本步骤构成；在潜空间推理中，这一过程发生在模型开口之前的内部。既可以用连续向量替换离散词元，也可以在后台生成真实文本词语，但不向用户显示。下面介绍其中一些方法。

\begin{itemize}
\item \vspace{0.5cm} 一个最简单的想法是使用“暂停词元”（pause tokens）。在标准自回归生成中，语言模型每输出一个词，就进行一步内部计算。有趣的是，研究发现，向模型输入占位词元，即暂停词元，就能在模型必须输出最终答案之前，给予它额外计算时间，在隐层中处理信息 \cite{goyal-etal:2024think}。这表明，大语言模型可以在不生成真实词语的情况下，进行无声计算或潜在思考。
\item \vspace{0.3cm} 另一种方式是通过知识蒸馏，将显式推理压缩到隐状态中。例如，隐式思维链（implicit chain of thought）方法训练学生模型，直接模仿更大教师模型在“出声思考”时的内部隐状态 \cite{deng-etal:2023implicit}，使学生学会跳过文本生成。大语言模型也可以通过自蒸馏从自身学习 \cite{shen-etal:2025codi}：强制潜在思考过程与基于文本的推理最终隐状态相匹配。在这两种方法中，潜在思考都体现为连续数学空间：模型不生成逐步可读文本，而是利用基于隐向量的表示计算推理步骤。
\item \vspace{0.3cm} 在连续潜空间中推理的另一种方式，是采用\mindex{连续思维链}（chain of continuous thought，\mindex{COCONUT}）范式 \cite{hao-etal:2025training}。在标准文本生成中，最后一层输出会转换为离散词元，再送回第一层，进行下一步计算。这迫使模型将内部计算压缩到单个词元中。COCONUT 在推理阶段绕过这一环节，直接将最后一层的输出向量送回输入层，作为\mindex{连续思维}（continuous thought）。反复执行后，模型在开始生成最终文本答案之前，先构建一条连续向量链。由于连续向量比离散词元能够容纳更多信息，模型可以同时探索多条推理路径，而无需过早确定某个具体词元。这一方式表明，大语言模型不必依赖人类语言来解决复杂问题；将推理过程保持在连续空间中，可以更高效地计算答案。
\item \vspace{0.3cm} 除了在连续空间中推理，也可以用普通文本词元表示内部思考。例如，在 Quiet-STaR 方法中 \cite{zelikman-etal:2024quietstar}，模型被训练为在读写普通文本的过程中静默思考。其思想很简单：预测下一个词元之前，先生成一小段隐藏的“思考”词元。可以将其视为用户看不到的内心独白或简短的心中草稿。生成这些隐藏思考后，模型借助它们预测真正的下一个词。如果这些思考有助于正确预测词语，模型就获得奖励，学到这种思考方式有用；如果思考无用或令人困惑，就学会将其丢弃。
\end{itemize}
\vspace{0.5cm}

\subsection{记忆与推理}

讨论大语言模型时，经常会出现一个问题：模型是真正在“思考”，还是仅仅记忆容量很大，正在复述训练数据中见过的类似答案？这引出了对模型中推理与记忆的比较。

简单来说，两者的区别可概括如下：

\begin{itemize}
\item \vspace{0.5cm} 记忆是向后看。模型调整参数，学习训练数据中的模式；面对类似输入时，只是重现学到的模式。
\item \vspace{0.3cm} 推理是向前走。模型不寻找现成答案，而是应用一般规则和逻辑，逐步构建思维链。推理的核心是泛化：即使面对训练集中从未出现的全新问题，也能借助逻辑工具解决。
\end{itemize}
\vspace{0.5cm}

在早期深度学习研究中，这条边界并不总是清晰。例如，\citet{zhang-etal:2021understand} 指出，深度神经网络具有巨大记忆容量，能够轻易记住标签完全随机的噪声数据。对大语言模型而言，庞大规模使其不仅能记住 Wikipedia，还能记住代码仓库中隐藏的细节。

这种记忆能力常会伪装成推理。例如，向模型提出一道经典算法题，只稍微修改数值，模型可能仍依赖记忆惯性，输出原答案，从而犯下基本错误。这通常称为“随机鹦鹉”（stochastic parrots）现象~\cite{bender-etal:2021dangers}，指模型只是盲目模仿、拼接记住的文本，而未真正理解背后的逻辑。

为了更清楚地分辨记忆与推理，可以考察两个不同维度。

\subsubsection{顿悟与泛化}

在探索模型如何从记忆转向真正理解时，研究者发现了一个有趣现象，称为“顿悟”（grokking）~\cite{power-etal:2022grokking}。他们训练一个小型 Transformer，学习模加法等抽象代数二元运算，并观察到十分有趣的现象：训练早期，模型在训练集上很快达到 100\% 准确率，但测试集准确率几乎为零。此时，模型显然只是在死记硬背等式。不过，如果不停止训练，让模型在远远超过过拟合发生的时刻后，仍反复学习这些数据，就会出现令人惊讶的变化：测试准确率突然在很短时间内飙升到 100\%，模型仿佛经历了一次“顿悟时刻”。


进一步研究表明，在“顿悟时刻”之前，模型参数以杂乱方式记住具体等式，例如 $2 + 3 = 5$~\cite{nanda-etal:2023progress}；但顿悟发生时，网络会在内部利用三角函数和傅里叶变换，形成通用数学表示。它找到了一个能够解释全部数据的通用算法。

这一实验说明，推理并非凭空出现。它往往在记忆遇到瓶颈时，或经过大量训练后，以高层次泛化状态涌现：网络组织其参数，找到最具普适性的解释。

\subsubsection{成本权衡}

从信息论和计算效率角度看，大语言模型会自然地根据成本平衡记忆与推理。一方面，记忆需要更大的参数容量。要记住世界上每个事实，模型规模就必须无限增长。另一方面，推理需要在模型运行时执行更多计算步骤。只要记住加法、乘法的一般规则，模型就能在生成过程中执行这些计算步骤，解决许多数学问题。

因此，大语言模型通常选择阻力最小的路径。对于常见事实，直接从参数或记忆中提取答案最快；对于罕见的多步任务，系统 2 思维则更合适。

这也解释了检索增强生成（retrieval-augmented generation，RAG）等技术为何流行。RAG 的实质，是将记忆负担转移给外部数据库，从而腾出参数空间用于推理。将知识检索与逻辑演绎分开后，模型无需记住不断变化的外部事实，而是知道如何查找和处理信息。

\subsection{大语言模型能真正推理吗？}

大语言模型能否真正推理，是当今人工智能领域争论最多的主题之一。看到模型完美解决数学问题，或编写复杂日程时，很容易相信它像人一样理解世界。不过，要回答这一问题，必须考察这些模型实际如何工作。

从本质上说，大语言模型通过预测下一个词运行。由于处理了海量文本，它们已完全掌握人类语言的规则、语法和结构，在认知科学中，这称为\mindex{形式语言能力}（formal linguistic competence）\cite{chomsky:1965aspects,mahowald-etal:2024dissociating}。它们准确知道词语应如何按逻辑排列，但在功能语言能力（functional linguistic competence）方面仍有困难，即运用语言理解和应对现实世界的能力 \cite{kambhampati-etal:2024llms}。人类通过身体与环境交互来学习，大语言模型却只知道文本的底层模式。由于没有植根于现实世界，它们缺乏对物理空间和严格因果关系的真正理解 \cite{mahowald-etal:2024dissociating}。

在规划任务上测试大语言模型时，这一局限十分明显 \cite{valmeekam-etal:2023planbench}。例如，在把桌上彩色积木移动成特定堆叠结构这类经典规划问题中，如果使用训练数据里常见的词，模型往往能够求解；但如果将积木名称和动作替换为无意义词，表现就会大幅下降 \cite{kambhampati-etal:2024can}。真正的推理系统不会在意名称，因为底层数学逻辑完全相同。模型之所以失败，是因为熟悉的文本模式被打破了。

此外，思维链提示等技术会迫使模型输出逐步过程，表现出某些缓慢、审慎的人类思考行为。但实际上，模型仍然只能严格向前推进，预测通常的逻辑解释应该是什么样子。如果模型犯错，或持有带偏差的假设，往往就会自信地生成有缺陷的解释，只为证明错误答案合理 \cite{chen-etal:2025reasoning}。

归根结底，要回答大语言模型能否“真正”推理，首先需要定义真正推理的含义。如果期望模型像人类一样思考，前面的讨论清楚表明，两者存在差别。但如果将推理纯粹视为功能过程，答案就有所不同。推理可以简单定义为：接收复杂问题，将其分解，应用逻辑规则，并得出正确结论的能力。模型逐步正确解出一个数学方程，或编写逻辑合理的代码处理数据时，就在执行一条有效逻辑链。即使这一能力来自大规模神经网络训练，而非类人认知，结果仍然是成功解决问题。随着模型规模增大、学习数据增多，它们会超越单纯记忆答案，开始捕捉逻辑背后的规则。

%----------------------------------------------------------------------------------------
% NEW SECTION
%----------------------------------------------------------------------------------------

\section{小结}

本章介绍了如何将大语言模型用于推理，这被视为 NLP 中最具挑战性的任务之一。重点讨论了在训练阶段和模型推理阶段提升推理能力的方法，还介绍了推理大语言模型的评估方法，以及一些尚待解决的问题。

大语言模型能够显式列出推理步骤，这一发现启发了新的研究方向和实际应用。当前研究不再只把推理看作回答孤立问题的手段，而是将其视为驱动更复杂、更自主行为的工具。

第一，推理能力与“长时间思考”密切相关，根本改变了大语言模型使用计算资源的方式。如第 \ref{sec:inference-time-reasoning} 节所述，给予模型更多思考时间和计算资源，就能解决更困难的问题。这使模型工作方式发生重大变化：现在不仅要在初始训练时分配算力，还需要弄清生成答案时如何最优分配算力。在系统层面解决这一问题，超出了本书范围。不过，构建高效、懂得该花多少时间思考的推理模型，是一个很有前景且实用的研究方向。

第二，推理能力是智能体推理的基础，使大语言模型从被动回答问题的机器，转为环境中的主动参与者。标准文本生成中的推理，是从提示到答案的直线过程；智能体推理则在思考、行动、观察的持续循环中运行 \cite{yao-etal:2023react}。然而，这一循环也带来重大挑战：模型很容易在循环中迷失或陷入困境。它们可能重复同样的错误、制定糟糕的计划，或在动作失败时，不知道如何从失败中恢复 \cite{shinn-etal:2023reflexion}。让模型长时间保持在正确轨道上，仍然非常困难。

第三，推理能力直接支持工具增强推理，克服大语言模型无法精确计算或获取实时事实等局限。可以训练模型认识自身限制，通过编写代码 \cite{gao-etal:2023pal} 或调用外部 API \cite{schick-etal:2024toolformer} 完成繁重工作，而不必强迫它猜测答案。不过，主要问题是教会模型何时需要工具，以及如何从多种工具中选择 \cite{qin-etal:2024toolllm}。此外，如果工具失效或返回异常数据，模型必须足够聪明，能够发现错误并再次尝试，而不是直接相信错误结果。

